% Use lualatex
\documentclass{llncs}
  
\usepackage[T1]{fontenc}
\usepackage[utf8]{inputenc}
\usepackage[russian,english]{babel}
\usepackage{fontspec}
\usepackage{lmodern}
\usepackage{mathtools}
%\usepackage{fullpage}
\usepackage{graphicx}
\usepackage{xspace}
\usepackage{tabularx}
\usepackage[lf]{ebgaramond}
\usepackage{biolinum}
\usepackage[cmintegrals,cmbraces]{newtxmath}
\usepackage{ebgaramond-maths}
\usepackage{multicol}
\usepackage{sectsty}
\usepackage[noend]{algpseudocode}
\usepackage{algorithm}
\usepackage{algorithmicx}
\usepackage[bottom]{footmisc}
\usepackage{caption}
%\captionsetup{font=small}
\captionsetup{labelfont={sf,bf}}
\usepackage{ragged2e}
\usepackage{xcolor}

\definecolor{codebg}{rgb}{.97,.97,.99}
\definecolor{gray}{rgb}{.4,.4,.4}
\definecolor{webgreen}{rgb}{0,.5,0}
\definecolor{Maroon}{cmyk}{0, 0.87, 0.68, 0.32}
\definecolor{RoyalBlue}{cmyk}{1, 0.50, 0, 0}
\definecolor{blood}{rgb}{.6,0,0}

\usepackage{listings}
\AtBeginDocument{%\counterwithin{lstlisting}{section}%
    \renewcommand\thelstlisting{\arabic{lstlisting}}}
\renewcommand{\lstlistingname}{Listing}
\DeclareCaptionFormat{listing}{#1#2#3}
\captionsetup[lstlisting]{format=listing,singlelinecheck=false}

\lstdefinestyle{zzpython}{%
    frame=ltb,
    framerule=0pt,
    aboveskip=3mm,
    belowskip=3mm,
    framextopmargin=.5mm,
    framexbottommargin=.5mm,
    framexleftmargin=1mm,
    framexrightmargin=1.5mm,
    showstringspaces=false,
    language=Python,
    columns=fullflexible,
    keepspaces=true,
    %basicstyle={\small\ttfamily},
    basicstyle={\fontsize{8pt}{9pt}\ttfamily},
    numbers=left,
    numberstyle={\fontsize{8pt}{9pt}\ttfamily\color{gray}},
    keywordstyle=\color{webgreen},
    commentstyle=\color{gray},
    stringstyle=\color{Maroon},
    identifierstyle=\color{RoyalBlue},
    breaklines=true,
    breakatwhitespace=true,
    tabsize=3,
    backgroundcolor=\color{codebg},
    captionpos=t,
}

\lstdefinestyle{zzc}{%
    frame=ltb,
    framerule=0pt,
    aboveskip=3mm,
    belowskip=3mm,
    framextopmargin=.5mm,
    framexbottommargin=.5mm,
    framexleftmargin=1mm,
    framexrightmargin=1.5mm,
    showstringspaces=false,
    language=C,
    columns=fullflexible,
    keepspaces=true,
    %basicstyle={\small\ttfamily},
    basicstyle={\fontsize{8pt}{9pt}\ttfamily},
    numbers=left,
    numberstyle={\fontsize{8pt}{9pt}\ttfamily\color{gray}},
    keywordstyle=\color{webgreen},
    commentstyle=\color{gray},
    stringstyle=\color{Maroon},
    identifierstyle=\color{RoyalBlue},
    breaklines=true,
    breakatwhitespace=true,
    tabsize=3,
    backgroundcolor=\color{codebg},
    captionpos=t,
}

\lstdefinelanguage{ARMAssembler}
    {%
        morekeywords={%
            adc,adcs,%
            add,adds,addw,%
            adr,%
            and,ands,%
            asr,asrs,%
            b,beq,bne,bcs,bcc,bmi,bpl,bvs,bvc,bhi,bls,bge,blt,bgt,ble,bal,bhs,blo,%
            bfc,%
            bfi,%
            bic,bics,%
            bkpt,%
            bl,%
            blx,%
            bx,%
            cbnz,cbz,%
            clrex,%
            clz,%
            cmn,%
            cmp,%
            cps,%
            cpy,%
            csdb,%
            dbg,%
            dmb,%
            dsb,%
            eor,eors,%
            isb,%
            it,itt,ite,ittt,itet,itte,itee,%
            itttt,ittte,ittet,ittee,itett,itete,iteet,iteee,%
            ldc,ldcl,ldc2,ldc2l,%
            ldm,ldmia,ldmfd,%
            ldmdb,ldmea,%
            ldr,%
            ldrb,%
            ldrbt,%
            ldrd,%
            ldrex,%
            ldrexb,%
            ldrexh,%
            ldrh,%
            ldrsb,%
            ldrsbt,%
            ldrsh,%
            ldrsht,%
            ldrt,%
            lsl,lsls,%
            lsr,lsrs,%
            mcr,lcr2,%
            mcrr,mcrr2,%
            mla,%
            mls,%
            mov,movs,movw,movt,%
            mrc,mrc2,%
            mrrc,mrrc2,%
            mrs,%
            msr,%
            mul,muls,%
            mvn,mvns,%
            neg,%
            nop,%
            orn,orns,%
            orr,orrs,%
            pkhbt,pkhtb,%
            pld,%
            pli,%
            pop,%
            pssbb,%
            push,%
            qadd,%
            qadd16,%
            qadd8,%
            qasx,%
            qdadd,%
            qdsub,%
            qsax,%
            qsub,%
            qsub16,%
            qsub8,%
            rbit,%
            rev,%
            rev16,%
            revsh,%
            ror,rors,%
            rrx,%
            rsb,rsbs,%
            sadd16,%
            sadd8,%
            sasx,%
            sbc,sbcs,%
            sbfx,%
            sdiv,%
            sel,sev,%
            shadd16,%
            shadd8,%
            shasx,%
            shsax,%
            shsub16,%
            shsub8,%
            smlabb,smlabt,smlatb,smlatt,%
            smlad,smladx,%
            smlal,%
            smlalbb,smlalbt,smlaltb,smlaltt,%
            smlald,smlaldx,%
            smlawb,smlawt,%
            smlsd,smlsdx,%
            smlsld,smlsldx,%
            smmla,smmlar,%
            smmls,smmlsr,%
            smmul,smmulr,%
            smuad,smuadx,%
            smulbb,smulbt,smultb,smultt,%
            smull,%
            smulwb,smulwt,%
            smusd,smusdx,%
            ssat,%
            ssat16,%
            ssax,%
            ssbb,%
            ssub16,%
            ssub8,%
            stc,stc2,%
            stm,stima,stmea,%
            stmdb,stmfd,%
            str,%
            strb,%
            strbt,%
            strd,%
            strex,%
            strexb,%
            strexh,%
            strh,%
            strht,%
            strt,%
            sub,subs,%
            svc,%
            sxtab,%
            sxtab16,%
            sxtah,%
            sxtb,%
            sxtb16,%
            sxth,%
            tbb,tbh,%
            teq,%
            tst,%
            uadd16,%
            uadd8,%
            uasx,%
            ubfx,%
            udf,%
            udiv,%
            uhadd16,%
            uhadd8,%
            uhasx,%
            uhsax,%
            uhsub16,%
            uhsub8,%
            umaal,%
            umlal,%
            umull,%
            uqadd16,%
            uqadd8,%
            uqasx,%
            uqsax,%
            uqsub16,%
            uqsub8,%
            usad8,%
            usada8,%
            usat,%
            usat16,%
            usax,%
            usub16,%
            usub8,%
            uxtab,%
            uxtab16,%
            uxtah,%
            uxtb,%
            uxtb16,%
            uxth,%
            vabs,%
            vadd,%
            vcmp,vcmpe,%
            vcvta,vcvtn,vcvtp,vcvtm,%
            vcvt,vcvtr,%
            vcvtb,vcvtt,%
            vdiv,%
            vfma,vfms,%
            vldm,%
            vldr,%
            vmaxnm,vminnm,%
            vmla,vmls,%
            vmov,%
            vmrs,%
            vmsr,%
            vneg,%
            vnmla,vnmls,vnmul,%
            vpop,%
            vpush,%
            vrinta,vrintn,vrintp,vrintm,%
            vrintx,%
            vrintz,vrintr,%
            vsel,%
            vsqrt,%
            vstm,%
            vstr,%
            vsub,%
            wfe,%
            wfi,%
            yield},%
        alsoletter={.,0,1,2,3,4,5,6,7,8,9},%
        alsodigit={?},%
        sensitive=false,%
        morestring=[b]",%
        morecomment=[s]{/*}{*/},%
        morecomment=[l]@,%
        morecomment=[l]//,%
    }[keywords,comments,strings]

\lstdefinestyle{zzarmasm}{%
    frame=ltb,
    framerule=0pt,
    aboveskip=3mm,
    belowskip=3mm,
    framextopmargin=.5mm,
    framexbottommargin=.5mm,
    framexleftmargin=1mm,
    framexrightmargin=1.5mm,
    showstringspaces=false,
    language=ARMAssembler,
    columns=fullflexible,
    basicstyle={\fontsize{8pt}{9pt}\ttfamily},
    basewidth=0.5em,
    numbers=left,
    numberstyle=\tiny\color{gray},
    keywordstyle=\color{webgreen},
    commentstyle=\color{gray},
    stringstyle=\color{Maroon},
    identifierstyle=\color{RoyalBlue},
    breaklines=true,
    breakatwhitespace=true,
    tabsize=3,
    backgroundcolor=\color{codebg}
}

\lstdefinestyle{zzc}{%
    frame=ltb,
    framerule=0pt,
    aboveskip=3mm,
    belowskip=3mm,
    framextopmargin=.5mm,
    framexbottommargin=.5mm,
    framexleftmargin=1mm,
    framexrightmargin=1.5mm,
    showstringspaces=false,
    language=C,
    columns=fullflexible,
    basicstyle={\fontsize{8pt}{9pt}\ttfamily},
    basewidth=0.5em,
    numbers=left,
    numberstyle=\tiny\color{gray},
    keywordstyle=\color{webgreen},
    commentstyle=\color{gray},
    stringstyle=\color{Maroon},
    identifierstyle=\color{RoyalBlue},
    breaklines=true,
    breakatwhitespace=true,
    tabsize=3,
    backgroundcolor=\color{codebg}
}

\newenvironment{forbidpagebreak}
    {\par\nobreak\vfil\penalty 0\vfilneg
     \vtop\bgroup}
    {\par\xdef\tpd{\the\prevdepth}\egroup
     \prevdepth=\tpd}

% llncs removes generation of PDF bookmarks, we must add them back
\usepackage{etoolbox}
\makeatletter
\let\llncs@addcontentsline\addcontentsline
\patchcmd{\maketitle}{\addcontentsline}{\llncs@addcontentsline}{}{}
\patchcmd{\maketitle}{\addcontentsline}{\llncs@addcontentsline}{}{}
\patchcmd{\maketitle}{\addcontentsline}{\llncs@addcontentsline}{}{}
\setcounter{tocdepth}{3}
\makeatother
\PassOptionsToPackage{hyphens}{url}
\usepackage[pdftex,bookmarks=true]{hyperref}
\hypersetup{
    bookmarksopen=true,
    bookmarksopenlevel=3,
    bookmarksnumbered=true,
    hidelinks,
    colorlinks,
    linkcolor={red!50!black},
    citecolor={green!50!black},
    urlcolor={blue!80!black}
}
%\hypersetup{hidelinks}
%\usepackage{attachfile2}

\makeatletter
  % Recover some math symbols that were masked by eb-garamond, but do not
  % have replacement definitions.
  \DeclareSymbolFont{ntxletters}{OML}{ntxmi}{m}{it}
  \SetSymbolFont{ntxletters}{bold}{OML}{ntxmi}{b}{it}
  \re@DeclareMathSymbol{\leftharpoonup}{\mathrel}{ntxletters}{"28}
  \re@DeclareMathSymbol{\leftharpoondown}{\mathrel}{ntxletters}{"29}
  \re@DeclareMathSymbol{\rightharpoonup}{\mathrel}{ntxletters}{"2A}
  \re@DeclareMathSymbol{\rightharpoondown}{\mathrel}{ntxletters}{"2B}
  \re@DeclareMathSymbol{\triangleleft}{\mathbin}{ntxletters}{"2F}
  \re@DeclareMathSymbol{\triangleright}{\mathbin}{ntxletters}{"2E}
  \re@DeclareMathSymbol{\partial}{\mathord}{ntxletters}{"40}
  \re@DeclareMathSymbol{\flat}{\mathord}{ntxletters}{"5B}
  \re@DeclareMathSymbol{\natural}{\mathord}{ntxletters}{"5C}
  \re@DeclareMathSymbol{\star}{\mathbin}{ntxletters}{"3F}
  \re@DeclareMathSymbol{\smile}{\mathrel}{ntxletters}{"5E}
  \re@DeclareMathSymbol{\frown}{\mathrel}{ntxletters}{"5F}
  \re@DeclareMathSymbol{\sharp}{\mathord}{ntxletters}{"5D}
  \re@DeclareMathAccent{\vec}{\mathord}{ntxletters}{"7E}

  % Change font for algorithm label.
  \renewcommand\ALG@name{\sffamily\bfseries Listing}
\makeatother

% Use the sans-serif fonts (for which bold is properly defined) for
% section headings. Also ensure that subsubsections get a number and
% that the number is displayed (to distinguish them from paragraphs).
\allsectionsfont{\sffamily}
\setcounter{secnumdepth}{3}
\pagestyle{plain}

\makeatletter
\renewenvironment{abstract}{%
      \list{}{\advance\topsep by0.35cm\relax\small
      \leftmargin=1cm
      \labelwidth=\z@
      \listparindent=\z@
      \itemindent\listparindent
      \rightmargin\leftmargin}\item[\hskip\labelsep
                                    \textsf{\textbf{\abstractname}}]}
    {\endlist}
\makeatother
\let\lemma\relax
\spnewtheorem{lemma}{Lemma}{\sffamily\bfseries}{\itshape}

\spnewtheorem{mtheorem}{Theorem}{\sffamily\bfseries}{\itshape}
\spnewtheorem*{mproof}{Proof}{\sffamily\bfseries\itshape}{\rmfamily}

\renewcommand{\algorithmicrequire}{\hbox to3.4em{\textbf{Input:}\hfill}}
\renewcommand{\algorithmicensure}{\hbox to3.4em{\textbf{Output:}\hfill}}

%\newcommand{\GF}{\mathrm{\textit{GF}}}
\newcommand{\GF}{GF}
\newcommand{\QR}{QR}
\newcommand{\bB}{\mathbb{B}}
\newcommand{\bF}{\mathbb{F}}
\newcommand{\bG}{\mathbb{G}}
\newcommand{\bN}{\mathbb{N}}
\newcommand{\bZ}{\mathbb{Z}}
\newcommand{\bR}{\mathbb{R}}
\newcommand{\cC}{\mathcal{C}}
\newcommand{\cE}{\mathcal{E}}
\newcommand{\cG}{\mathcal{G}}
\newcommand{\rM}{\text{M}}
\newcommand{\rS}{\text{S}}
\newcommand{\rmb}{\text{m}_\text{b}}
\newcommand{\neutral}{\mathbb{O}}
\newcommand{\vol}{\text{\textsf{vol}}}
\newcommand{\bitlength}{\text{\textsf{len}}}
\newcommand{\MAC}{\text{\textsf{MAC}}}
\newcommand{\MAClen}{\text{\textsf{MAClen}}}
\newcommand{\cc}{\text{\textsf{enc}}}
\newcommand{\cclen}{\text{\textsf{clen}}}
\newcommand{\Setup}{\text{\textsf{Setup}}}
\newcommand{\Eval}{\text{\textsf{Eval}}}
\newcommand{\Verify}{\text{\textsf{Verify}}}
\newcommand{\VDF}{\text{\textsf{VDF}}}
\newcommand{\VDFlen}{\text{\textsf{VDFlen}}}
%\newcommand{\smod}[1]{\,\,\,(\text{mod}^{\pm} #1)}
%\newcommand{\smodnospace}[1]{(\text{mod}^{\pm} #1)}
\newcommand{\sign}{\text{\textsf{sign}}}
\newcommand{\trace}{\text{Tr}}
\newcommand{\qsolve}{\text{QSolve}}
\newcommand{\rev}{\text{rev}}

\newcommand{\ezut}{\textsc{ezut}\xspace}
\newcommand{\xwj}{\textsc{xwj}\xspace}

\newcommand{\eqvm}[3]{#1 \equiv #2 \,[#3]}
\newcommand{\umod}[2]{#1 \,\,\text{mod}^{+}\, #2}
\newcommand{\smod}[2]{#1 \,\,\text{mod}^{\pm}\, #2}
\newcommand{\mont}[1]{\mathcal{M}(#1)}

\raggedbottom

\begin{document}

\title{\textsf{Montgomery Multiplication\\in Signed Redundant Representations}}

\author{Thomas Pornin}
\institute{NCC Group, \email{thomas.pornin@nccgroup.com}}

\maketitle
\noindent\makebox[\textwidth]{19 April, 2026}
%FIXME: adjust date when (re)publishing

\begin{abstract}
In this paper, we explore the use of Montgomery multiplication with a
multi-limb redundant representation of integers, in particular in
combination with signed reduction factors. We develop techniques that
are particularly suited to software platforms on which carry propagation
is expensive, in particular RISC-V CPUs which lack hardware support for
carries. We also show how to perform a whole-primitive range analysis
that demonstrates that overflows are not possible, thus allowing liberal
use of unreduced limb-wise additions and subtractions, which are small
and fast. The implementation and analysis techniques are illustrated in
a codegolfing exercise, to produce size-optimized implementations of
ECDSA signature verification over NIST curve P-256; use of a virtual CPU
with a custom instruction set with byte-size encoding (``bytecode'') allows
the production of an implementation as small as 848 bytes on x86 CPUs (in
64-bit mode); RISC-V (984 bytes), Armv8-A (1136 bytes) and portable C
implementations (about 2200 to 2800 bytes) are also provided. In the
process, an AI is utterly discomfited.
\end{abstract}

% ----------------------------------------------------------------------

\section{Introduction}\label{sec:intro}

Computations over modular integers are pervasive to many cryptographic
schemes, including classic algorithms such as RSA\cite{RivShaAdl1978}
and Diffie-Hellman\cite{DifHel1976}, elliptic curve signature
schemes (ECDSA\cite{Fips186,X962}, EdDSA\cite{BerDuiLanSchYan2012}) and key
echange methods (ECDH\cite{X963}), pairing-based cryptography over
pairing-friendly curves, but also lattice-based algorithms such as
ML-KEM\cite{Fips203}, ML-DSA\cite{Fips204}, Falcon\cite{FalconRound3}
(soon-to-be FN-DSA) and schemes relying on isogenies between
supersingular curves. There is an ample literature on the subject; a
most popular method for implementing arithmetic operations modulo some
(odd) integer is the use of Montgomery multiplication\cite{Mon1985}.

This work is, at its core, an exploration of some variations on the use
of \emph{redundant} big integer representations with Montgomery
multiplication: values are split into limbs which are strictly shorter
than the machine word size, and they are allowed to extend beyond their
nominal range. Redundant representations have long been studied, even
before computers were invented\cite{Col1726}, for efficient
computations\cite{Bri1983,Hon1989} and in particular for bounded carry
propagation techniques in custom hardware (ASIC)\cite{PhaKor1994}.
Redundant representations with signed limbs can efficiently support
operations modulo some special-format moduli\cite{Ber2006}. The main
potential advantage is the ability to perform additions and subtractions
in a limb-wise way, without any reduction; moreover, the extra limb
range can be leveraged to perform fewer carry propagations, which are
detrimental to out-of-order and parallel execution, and also lack
hardware support in some platforms (e.g. RISC-V). We furthermore
consider use of \emph{signed} values for the Montgomery multiplicative
factors, as well as signed or unsigned normalization. These variants are
described in section~\ref{sec:montymul}.

To get a concrete context for the discussion, we implemented ECDSA
signature verification over the NIST curve P-256 with optimization for
code size rather than speed; this is a somewhat uncommon goal, which has
some niche usages (e.g. boot-time verification of a signature on
Flashable firmware, the verification code must fit in a small ROM). This
was the occasion of a competition in which a developer (myself) was
pitted against an LLM-based AI (Claude Code) over the course of about 8
days, with Intel x86 architecture as target (in 64-bit mode). The long
history of the x86 platform made its instruction set especially rich in
abstruse size optimizations, for which the pattern-matching behaviour of
the LLM should be effective, especially since x86 assembly writing is a
subject for which the Internet provides an extremely large corpus of
examples. The bot was ultimately defeated, though it did provide some
tough opposition; it also got it wrong at some point and produced code
which deviated from the ECDSA specification. This is described in
section~\ref{sec:showdown}. The outcome is instructive about what can
and cannot be expected from LLMs, and how they may fail in the context
of implementing cryptographic primitives.

In section~\ref{sec:bytecode}, my own implementations are described, not
only the x86-optimized code, but also 64-bit Arm (Armv8-A architecture)
and 64-bit RISC-V (RV64 with I, M and C extensions), as well as C code
with 64-bit and 32-bit words. The assembly implementations use a custom
dedicated virtual machine with a few instructions that fit in a single
byte each (``bytecode''), allowing to pack all the needed operations for
ECDSA in a compact format.

Use of limb-wise additions and subtractions with no reduction implies
that if too many such operations are performed in sequence, with no
interleaved multiplications, then limb values may overflow. An
implementation using that method can be deemed correct only if it can be
proven never to reach an overflow condition. This requires an extensive
and tight enough range analysis. In section~\ref{sec:range}, the
bytecode from the previous section is leveraged into a Python-based
interpreter which makes a comprehensive range analysis for all values
throughout the entire ECDSA operation. This framework can moreover be
reused over other word sizes, radices, and number of limbs; this is how
my implementations with 32-bit words were produced.

All the implementations described here, as well as the range analysis
framework, are available at:
\begin{center}
    \url{https://github.com/pornin/small-ecdsa}
\end{center}

% ----------------------------------------------------------------------

\newpage
\section{Montgomery Multiplication and Variants}\label{sec:montymul}

We describe here the classic Montgomery multiplication, and how it can
be used in a redundant representation with signed limbs. For a
well-known description of the classic algorithms, see the famous
Handbook of Applied Cryptography (\cite{MenOorVan1996}, section 14.3).
Of all the discussion presented here, about none is really novel, except
possibly the use of signed Montgomery factors $f_i$, and the
pre-shifting of operands to ease the low/high split in long
multiplication results.

\subsection{Notations}

For an integer $m > 0$, and two integers $x$ and $y$, we write
``$\eqvm{x}{y}{m}$'' to express that $m$ divides $x - y$ without
prescribing a particular range on either $x$ or $y$. For modular
\emph{unsigned reduction}, we write ``$\umod{x}{m}$'' to designate the
unique integer $y$ such that $\eqvm{x}{y}{m}$ and $y \in [0, m-1]$; for
\emph{signed reduction}, ``$\smod{x}{m}$'' is the unique integer $y$
such that $\eqvm{x}{y}{m}$ and $-m/2 \leq y < +m/2$. Note in particular
that when $m$ is even, the lower end $-m/2$ is included in the signed
reduction output, but the upper end $+m/2$ is not: this corresponds to
how usual software platforms interpret signed words with two's
complement notation, but it differs from the ``$\text{mod}^\pm$''
notation used in the ML-DSA scheme.

By extension, we may use these notations with \emph{rationals}: for
integers $x$, $y$ and $z$ such that $y$ is coprime with $m$, we write
$\eqvm{x/y}{z}{m}$ to signify that $\eqvm{x}{yz}{m}$; values
$\umod{x/y}{m}$ and $\smod{x/y}{m}$ denote integers $z$ such that
$\eqvm{x/y}{z}{m}$, with $0 \leq z < m$ and $-m/2\leq z < +m/2$,
respectively.

\subsection{Montgomery Multiplication and Representation}%
\label{sec:montymul:repr}

Let $m > 0$ be an integer; we want to perform arithmetic operations with
integers modulo $m$. Let $R > m$ be some constant which is relatively
prime to $m$, and such that Euclidean division by $R$ on integers is
easy. In all of the following, $m$ is an odd integer, and $R$ is a power
of $2$, so that the quotient and remainder of a division by $R$ can be
obtained with a right-shift and a truncation on the binary
representation, respectively; however, the concept can be used with
other values of $R$, as long as Euclidean division by $R$ is efficient.

\emph{Montgomery reduction} of an integer $z$ consists in computing $f =
\umod{-z/m}{R}$, then $z' = (z + fm)/R$. The computation of $f$ (the
``Montgomery factor'') is done by multiplying $z$ with the value $C =
\umod{-1/m}{R}$, which depends only on the modulus $m$ and is assumed to
be precomputed ($C$ exists because $m$ is relatively prime to $R$); this
multiplication is done modulo $R$, which is efficient (by definition of
$R$). This specific factor $f$ implies that $\eqvm{z + fm}{0}{R}$; the
division $(z + fm)/R$ is therefore exact: the output $z'$ is an integer.

If we consider values modulo $m$, then $\eqvm{z'}{z/R}{m}$ (adding a
multiple of $m$ does not change the value modulo $m$). The point of
Montgomery reduction is seen by considering ranges: since $0\leq f < R$,
we obtain that $|z'| < |z/R| + m$; in particular, if $0\leq z < Rm$,
then $0\leq z' < 2m$.

\emph{Montgomery multiplication} is the combination of plain integer
multiplication and Montgomery reduction. Given two operands $x$ and $y$,
Montgomery multiplication computes $z = xy$ then applies Montgomery reduction
and returns $z' = (xy + fm)/R$. If $x$ and $y$ are both in $[0, m-1]$, then
$0\leq z' < 2m$ and can be fully reduced modulo $m$ with a single conditional
subtraction of $m$ (subtract $m$ from $z'$ if and only if $z' \geq m$).
We denote with ``$*$'' the Montgomery multiplication; for inputs $x$ and
$y$, we thus obtain that $\eqvm{x*y}{xy/R}{m}$.

\emph{Montgomery representation} of an integer $x$ is another integer
$\mont{x}$ such that $\eqvm{\mont{x}}{xR}{m}$. Montgomery representation
is compatible with addition, and with multiplication, provided that we
use Montgomery multiplication: for all integers $x$ and $y$,
$\eqvm{\mont{x} + \mont{y}}{\mont{x + y}}{m}$, and
$\eqvm{\mont{x} * \mont{y}}{\mont{xy}}{m}$. This finally yields the
classic implementation method based on Montgomery representation:
\begin{itemize}

    \item All values are kept in Montgomery representation, normalized
    to $[0, m-1]$. Conversion to Montgomery representation from a freshly
    decoded $x$ integer is done by computing
    $x * (\umod{R^2}{m}) = \mont{x}$, the constant $\umod{R^2}{m}$ being
    precomputed.

    \item Additions are done with a plain integer addition, followed by
    a conditional subtraction of the modulus $m$ to ensure a result in
    $[0, m-1]$. Similarly, subtractions are done over plain integers, then
    adding back $m$ if the result is negative.

    \item Multiplications are performed with Montgomery multiplication,
    which yields the correct result since we work with Montgomery
    representations. A conditional subtraction of the modulus $m$ is
    applied at the end of the Montgomery multiplication to ensure that
    the output is in $[0, m-1]$.

    \item On output, right before encoding results, values are converted
    back from Montgomery representation to normal representation by
    applying Montgomery reduction. This can be implemented with a
    Montgomery multiplication: $\eqvm{x}{\mont{x} * 1}{m}$.

\end{itemize}
One may note that if $x\in [0, m-1]$ then Montgomery reduction of $x$ is
also in $[0, m-1]$: we have $0\leq f\leq R-1$, hence $0\leq x + fm\leq m
- 1 + (R - 1)m = Rm - 1$; therefore, the final division by $R$
necessarily yields a value strictly lower than $m$. This means that the
final Montgomery reduction does not need the conditional subtraction of
the modulus to ensure a canonical output. Notwithstanding, such
operations are usually not time-critical, and it is easier to simply use
the Montgomery multiplication routine with the constant $1$, even though
that implies some useless operations (plain multiplication by $1$ is a
no-operation).

\subsection{Multi-Limb Montgomery Multiplication}

For some cryptographic schemes (typically lattice-based schemes), all
involved integers fit in a machine word and the implementation can
follow the formulas described above. However, some others need to work
with larger integers, which must then be split over several words. Let
$w$ be the machine word size, in bits; the implementation platform can
compute additions and subtractions over values by working, implicitly,
modulo $2^w$. Each $w$-bit word can be considered as either
\emph{signed} or \emph{unsigned}:
\begin{itemize}

    \item In signed representation, the word contains a value in
    $[-2^{w-1}, 2^{w-1}-1]$ (with two's complement notation for
    negative values).

    \item In unsigned representation, the word contains a value in
    $[0, 2^w - 1]$.

\end{itemize}
A given pattern of $w$ bits can be converted between the two
representations by implicit reinterpreting, which usually has no reality
at the assembly level: the CPU does not care whether you think of some
values as signed or not, it changes nothing to the additions,
subtractions, and multiplications modulo $2^w$. Mathematically, if we
number bits from $0$ to $w-1$ (in low-to-high order), then unsigned
representation means that bit $w-1$ has weight $2^{w-1}$, but weight
$-2^{w-1}$ if the value is signed.

``Long multiplications'' are multiplications with an output that spans
over two words (that is, multiplications which are done over plain
integers, \emph{not} modulo $2^w$). For these operations, signedness
matters; thus, CPUs usually offer two instructions, one for
multiplication of signed values, and one for multiplication of unsigned
values\footnote{RISC-V CPUs also have a third kind, for mixed
signed/unsigned multiplications.}. For input integers $a$ and $b$, each
being the signed or unsigned interpretation of a $w$-bit word, long
multiplication returns two words, the ``low word'' and the ``high
word''. The low word is to be interpreted as an unsigned value $c_L \in
[0, 2^w - 1]$; the high word is interpreted as value $c_H$ with unsigned
representation if both $a$ and $b$ are unsigned, or signed
representation otherwise. We then have $ab = c_L + 2^w c_H$.

These operations map to the abilities of usual software platforms, at
least at the hardware level. In programming languages, signed and
unsigned interpretations are usually distinguished with types, but
explicit conversions (e.g. with Rust's ``\verb+x as u32+'' syntax, or
Go's ``\verb+uint32(x)+'') can be used to move between the two
interpretations at no cost at the hardware level.

In the following, we will use a Python-based pseudocode notation with
arithmetic operations on words specified with dedicated functions; in
the descriptions below, we use $x_\text{s}$ and $x_\text{u}$ to denoted
the contents of the word \verb+x+ interpreted as signed and unsigned,
respectively (i.e. $-2^{w-1} \leq x_\text{s} < 2^{w-1} - 1$,
$0\leq x_\text{u} < 2^w$, $x_\text{s} = \smod{x_\text{u}}{2^w}$
and $x_\text{u} = \umod{x_\text{s}}{2^w}$).

\newcommand{\makecodedesclabel}[1]{\texttt{\small #1}\hfill}
\newenvironment{codedesc}
    {\begin{list}
        {}
        {\setlength{\labelwidth}{10.5em}
         \setlength{\labelsep}{0.5em}
         \setlength{\itemsep}{1ex}
         \setlength{\leftmargin}{11em}
         \setlength{\rightmargin}{0em}
         \let\makelabel=\makecodedesclabel
        }
    }
    {\end{list}}

\begin{codedesc}
    \item[add(a, b)] Addition: returns $\umod{a_\text{u} + b_\text{u}}{2^w}$.

    \item[add\_carry(a, b, c\_in)] Addition with carry: this computes
    $d = a_\text{u} + b_\text{u} + 1$ if $c_{\text{in}}$ is \verb+True+,
    or $d = a_\text{u} + b_\text{u}$ if $c_{\text{in}}$ is \verb+False+.
    Two values are returned: $\umod{d}{2^w}$, and the output carry
    $c_{\text{out}}$ which is \verb+True+ if $d \geq 2^w$, or
    \verb+False+ if $d < 2^w$.

    \item[sub(a, b)] Subtraction: returns $\umod{a_\text{u} - b_\text{u}}{2^w}$.

    \item[sub\_borrow(a, b, c\_in)] Subtraction with borrow: this computes
    $d = a_\text{u} - b_\text{u} - 1$ if $c_{\text{in}}$ is \verb+True+,
    or $d = a_\text{u} - b_\text{u}$ if $c_{\text{in}}$ is \verb+False+.
    Two values are returned: $\umod{d}{2^w}$, and the output borrow
    $c_{\text{out}}$ which is \verb+True+ if $d < 0$, or
    \verb+False+ if $d \geq 0$.

    \item[lsh(a, n)] Left-shift: returns $\umod{2^n a_\text{u}}{2^w}$. Shift
    count \verb+n+ must be a non-negative constant integer. At the binary
    level, bits are shifted to the left (toward more significant places),
    and zeros are injected in the vacated low positions.

    \item[rsh\_s(a, n)] Right-shift (arithmetic): returns
    $\umod{\lfloor a_\text{s}/2^n\rfloor}{2^w}$. Shift count \verb+n+ must
    be a non-negative constant integer. At the binary level, bits are
    shifted to the right (toward less significant places) and copies of the
    sign bit (leftmost bit) of the input are injected in the vacated high
    positions.

    \item[rsh\_u(a, n)] Right-shift (logical): returns
    $\umod{\lfloor a_\text{u}/2^n\rfloor}{2^w}$. Shift count \verb+n+ must
    be a non-negative constant integer. At the binary level, bits are
    shifted to the right (toward less significant places) and zeros are
    injected in the vacated high positions.

    \item[mul(a, b)] Multiplication: returns
    $\umod{a_\text{u} b_\text{u}}{2^w}$.

    \item[longmul\_s(a, b)] Long multiplication (signed): this computes
    $d = a_\text{s} b_\text{s}$, then returns that value as two words:
    the low word $\umod{d}{2^w}$ and the high word
    $\umod{\lfloor d/2^w\rfloor}{2^w}$. The inputs are interpreted
    as signed values for the purpose of computing the product $d$.

    \item[longmul\_u(a, b)] Long multiplication (unsigned): this computes
    $d = a_\text{u} b_\text{u}$, then returns that value as two words:
    the low word $\umod{d}{2^w}$ and the high word
    $\umod{\lfloor d/2^w\rfloor}{2^w}$. The inputs are interpreted
    as unsigned values for the purpose of computing the product $d$.

    \item[is\_zero(x)] Comparison against zero: returns \verb+True+ if
    $x_\text{u} = 0$, or \verb+False+ otherwise.

\end{codedesc}

These functions correspond to hardware instructions offered by common
CPUs, with some variations, in particular how carries and borrows are
handled:
\begin{itemize}
    \item On x86 and Arm CPUs, there is a single carry flag, stored in
    the flags register, which all instructions must share. Recent x86
    support extra instructions (\verb+adcx+ and \verb+adox+) which allow
    maintaining two interleaved addition chains with separate carries,
    though there is no corresponding support for
    subtractions.\footnote{It should be understood that on modern
    \emph{large} CPUs, with out-of-order execution, the physical CPU
    contains many separate carry flags to which the programming-level
    unique flag is dynamically mapped. Having a single carry still
    implies some constraints on both the compiler or assembly developer,
    and on how much parallelism the CPU can infer from the code.}

    \item On Arm CPUs specifically, the CPU carry flag is reversed for
    borrows: the carry flag is set when there is no borrow, and is cleared
    when there is a borrow. This is an implementation detail which is not
    relevant to Montgomery multiplication, since we do not mix carry flags
    between additions and subtractions. In this paper, we treat carries
    and borrows as Boolean values \verb+True+ or \verb+False+.

    \item PowerPC CPUs (and later POWER CPUs) offer eight separate carry
    flags. It should be noted that some PowerPC implementations (especially
    the ``PPE'' core found in the PlayStation 3 and Xbox 360 consoles)
    handle only the first flag natively, and resort to slow microcode for
    the other seven flags.

    \item As opposed to PowerPCs' great trove of carry flags, RISC-V
    CPUs do not have a carry flag at all. For these CPUs,
    \verb+add_carry()+ requires no fewer than \emph{five} instructions
    in all generality (but only two instructions if either the input
    carry is known to be \verb+False+, or the output carry is ignored).
    The inefficiency of carry propagation on RISC-V was the initial
    motivation for this work.

    \item SIMD instructions (SSE2/AVX2/AVX-512 on x86, NEON on Arm, etc)
    typically do not have carry flags. Implementation using these
    instructions must resort to solutions similar to those for RISC-V
    code.

\end{itemize}

\begin{figure}
\begin{forbidpagebreak}
\lstinputlisting[caption={\label{alg:mmul-classic}Multi-limb modular arithmetics (classic representation)},style=zzpython]{mmul-classic.py}
\end{forbidpagebreak}
\end{figure}

With these notations, the pseudocode in listing~\ref{alg:mmul-classic}
shows how a classic multi-limb modular integer implementation would
compute additions, subtractions and Montgomery multiplications. Each
value $x$ is represented in radix $2^w$ as an array of $k$ words
$(x_i)_{0\leq i < k}$, which are implicitly considered unsigned (in
$[0, 2^w - 1]$). The words $(x_i)$ encode $x$ with the rule that
$x = \sum_{i=0}^{k-1} x_i 2^{iw}$. Modular integers are kept in
canonical unsigned format: inputs and outputs of modular addition,
subtraction and Montgomery reduction are always in $[0, m - 1]$, with
$m$ being the modulus.

The internal functions \verb+bigint_add_carry()+ and
\verb+bigint_sub_carry()+ compute multi-limb additions and subtractions
on non-modular integers, and return the result modulo $2^{kw}$ as well
as an output carry or borrow; they are used for modular addition
(\verb+mod_add()+) and subtraction (\verb+mod_sub()+) to ensure that the
output is canonical. The internal routine \verb+bigint_add_mul()+ is
used for Montgomery multiplication; it adds $gx$ to a mutable value $d$,
with $g$ being a single word, $x$ an integer over $k$ words, and $d$
being ``large enough'', meaning that it has at least $k+1$ words,
possibly $k+2$ (it can be shown that within the Montgomery
multiplication as implemented here, $d$ is always large enough).
Montgomery multiplication itself is \verb+mod_montymul()+ and uses $R =
2^{kw}$. Implicitly used parameters are the number of limbs $k$, the
modulus $m$, and the Montgomery coefficient $\umod{-1/m}{2^w}$ (denoted
\verb+m0i+ in the code). Since the inputs are canonical, the core of the
Montgomery multiplication always returns a value in $[0, 2m - 1]$ and
a single conditional subtraction of the modulus (lines 55-56 in the code)
is sufficient to ensure a canonical output.

Not shown in listing~\ref{alg:mmul-classic} are the decoding function
(split of the input into $k$ words of $w$ bits), conversion to
Montgomery representation (done with \verb+mod_montymul()+ and the
precomputed constant $\umod{R^2}{m}$), Montgomery reduction
(\verb+mod_montymul()+ with $1$ as second operand), and reassembly of
the words into a suitable output format, e.g. a sequence of bytes.

Some variants are possible; in particular, the implementation might not
insist on a strictly canonical output, and may instead work internally
with a somewhat larger range, e.g. $[0, R-1]$ (all sequences of $k$
words being then possible). Code working with a fixed compile-time
modulus, as is typical of optimized elliptic curve implementations, may
choose to unroll loops and use statically assigned registers rather than
arrays for values; the routines may also be subject to different degrees
of inlining. \emph{A contrario}, for a large modulus, and especially for
a modulus whose contents and even size are known only at runtime, actual
arrays and indexing are likely to be used, and the two
\verb+bigint_add_mul()+ routines called in each iteration of the
\verb+mod_montymul()+ loop may be interleaved so as to minimize such
memory traffic.

The salient point of the classic implementation is its reliance on carry
propagation, both at the hardware level (the \verb+add_carry()+ and
\verb+sub_borrow()+ primitives) and through the propagation of the
``\verb+delayed+'' word on \verb+bigint_add_mul()+. Modular addition and
subtraction both involve two chains of instructions with linear
dependencies (along each chain, the carry/borrow output of an
instruction is used as input for the next), which reduces parallelism
opportunities, especially when using SIMD instructions or out-of-order
execution; on platforms without support for carries (such as RISC-V),
these operations are particularly expensive, both in code size and in
execution speed. Additions in \verb+bigint_add_mul()+ also represent
long dependency chains (each value of the \verb+delayed+ variable
depends on the output of the previous iteration, which itself depends on
the previous value of \verb+delayed+). The conditional choices
(``\verb+if+'' clauses in \verb+mod_add()+, \verb+mod_sub()+ and
\verb+mod_montymul()+, on lines 38, 44 and 56, respectively) may
moreover involve an extra value scan to support a constant-time
implementation, in case the involved values are secret. Our improved
implementation method with a redundant representation was designed in
particular to avoid most of these carry propagation issues.

\subsection{Redundant Representation}\label{sec:montymul:redundant}

A \emph{redundant representation} of a big integer uses a radix which is
strictly smaller than the word size. Here, we will use radix $2^s$ for
some integer $s < w$ (with $w$ being the word size). Redundancy comes
from the idea that individual limbs may have values beyond the nominal
$[0, 2^s - 1]$ range, which implies that a given integer may have
several possible representations. Redundant representations are a
well-known area (see section~\ref{sec:intro} for references). In the
version described here, we are investigating the effect of some variants
that involve use of signed values in various places.

In all generality, the mathematical description is the following. For
a given integer $x$, we use a representation over $k$ limbs in radix
$2^s$:
\begin{equation*}
    x = \sum_{i=0}^{k-1} x_i 2^{si}
\end{equation*}
The limbs are, at this level, plain integers which are not limited to the
$[0, 2^s - 1]$ range; they can be larger than $2^s - 1$, or negative. The
representation of $x$ as $(x_i)$ is not unique. We can do additions and
subtractions in a limb-wise way:
\begin{equation*}
    x + y = \sum_{i=0}^{k-1} (x_i + y_i) 2^{si}
\end{equation*}
In a redundant representation, we can move some bits from a limb to the
next, using:
\begin{equation*}
    x_i + x_{i+1} 2^s = (\umod{x_i}{2^s}) + (x_{i+1} + \lfloor x_i/2^s \rfloor) 2^s
\end{equation*}
In plain words, we can split a limb into its low $s$ bits (a value in
$[0, 2^s - 1]$) and its upper bits, which can be added to the next limb
with a right-shift by $s$ bits. Since limb values are potentially
negative, we of course consider an arithmetic right shift here. A signed
analog exists, in which we force the lower limb into $[-2^{s-1}, 2^{s-1}
- 1]$:
\begin{equation*}
    x_i + x_{i+1} 2^s = (\smod{x_i}{2^s}) + \left(x_{i+1} + \frac{x_i - (\smod{x_i}{2^s})}{2^s}\right) 2^s
\end{equation*}
This analog will be considered in an implementation variant (``signed
normalization'').

Montgomery multiplication uses the $s$-bit Montgomery coefficient
$g = \umod{-1/m}{2^s}$, which depends on the modulus $m$ but not on the
operands $x$ and $y$. The process starts with an integer $t_0$, as a
$k$-limb sequence $(t_{0,j})$, with all limbs equal to zero. Then, for
$i = 0$ to $k-1$, we define:
\begin{eqnarray*}
    u_{i,j} &=& t_{i,j} + (\umod{x_i y_j}{2^s}) + \lfloor x_i y_{j-1} / 2^s \rfloor \\
    f_i &\equiv& g u_{i,0} \,\,[2^s] \\
    v_{i,j} &=& u_{i,j} + (\umod{f_i m_j}{2^s}) + \lfloor f_i m_{j-1} / 2^s \rfloor \\
    t_{i+1,0} &=& v_{i,1} + v_{i,0}/2^s \\
    t_{i+1,j} &=& v_{i,j+1} \ \ \ (\text{for\ }j\geq 1)
\end{eqnarray*}
These equations are to be understood with $j$ ranging from $0$ to $k$
(inclusive), with the convention that limbs with an out-of-bounds index
(e.g. $m_{-1}$ or $x_k$) have value zero. Intermediate values $u_i$ and
$v_i$ expand over $k+1$ limbs. The values $f_i$ are the Montgomery
factors; note that they are defined above only as a congruence class
modulo $2^s$: the operation works regardless of which modular reduction
is used. Classic Montgomery reduction would set $f_i = \umod{g
u_{i,0}}{2^s}$ (factor in $[0, 2^s - 1]$), but here we will use $f_i =
\smod{g u_{i,0}}{2^s}$ (factor in $[-2^{s-1}, 2^{s-1} - 1]$). The core
idea of the Montgomery multiplication still applies: if we work modulo
$2^s$, then:
\begin{eqnarray*}
    v_{i,0} &\equiv& u_{i,0} + f_i m \,\,[2^s] \\
            &\equiv& u_{i,0} + m (-1/m) u_{i,0} \,\,[2^s] \\
            &\equiv& 0 \,\,[2^s]
\end{eqnarray*}
so that the division of $v_{i,0}$ by $2^s$ (in the definition of
$t_{i+1,0}$) is exact, and the dropped value $\umod{v_{i,0}}{2^s}$ is in
fact zero. This implies that $t_{i+1} = (t_i + x_i y + f_i m)/2^s$, and
the $k$ iterations collectively compute $t_k = (xy + fm)/2^{ks}$ for an
integer $f$ whose $k$-limb representation is $(f_i)$. This follows the
classic Montgomery multiplication for the values, with $R = 2^{ks}$,
though most carry and upper bits propagations have been removed (only
the upper bits of $v_{i,0}$ are propagated). Since the output limbs (in
$t_k$) tend to be quite large, we follow Montgomery multiplication with
a single upper bits propagation pass, which we will later call
``normalization''.

Listing~\ref{alg:mmul-redundant} shows a pseudocode implementation of
modular addition, subtraction and Montgomery multiplication in a redundant
representation. We will now highlight some optimizations and other
implementation variants.

\begin{figure}
\begin{forbidpagebreak}
\lstinputlisting[caption={\label{alg:mmul-redundant}Multi-limb Montgomery arithmetics (redundant representation)},style=zzpython]{mmul-redundant.py}
\end{forbidpagebreak}
\end{figure}

\paragraph{Limb-wise additions and subtractions.}
The functions for addition (\verb+mod_add()+) and subtraction
(\verb+mod_sub()+) operate in a limb-wise way, with no carry/borrow
propagation; in fact, listing~\ref{alg:mmul-redundant} uses no
\verb+add_carry+ or \verb+sub_borrow+ at all, which improves speed and
reduces code size on platforms that don't have hardware support for
carries. The limb additions and subtractions are also independent of
each other and can be carried in parallel, which may efficiently
leverage SIMD units such as SSE2. The flip side of this parallelism is
that limbs grow beyond the nominal range; if too many such operations
are performed in a row, limbs may overflow the representable range in
$w$-bit words. Most of the range analysis in section~\ref{sec:range} is
about proving that such overflows cannot happen in the code that uses
these redundant integers.

Apart from limb growth, unreduced additions and subtractions also entail
value growth: if the two source operands live in ranges $[a_L, a_H]$ and
$[b_L, b_H]$, then their limb-wise sum is in $[a_L + b_L, a_H + b_H]$.
Excessive value growth is \emph{also} a concern, and calls for
intermediate reduction steps that help prevent it. As will be detailed
below, limb growth can be controlled with normalization, and Montgomery
multiplication (which includes a normalization step) reduces both limb
and overall value ranges.

\paragraph{Signed limbs.}
Limb values are interpreted as signed integers; this is apparent in the
use of \verb+longmul_s+ in the Montgomery multiplications (line 21), and
an arithmetic right-shift in normalization (line 33). Support of
negative limb values allows a simple limb-wise subtraction; if limbs had
to be kept non-negative at all times, then subtraction would involve
either carry propagation, or at least adding extra multiples of the
modulus, which would entail further limb growth, as well as more
instructions to execute.

\paragraph{Normalization.}
To contain limb growth, normalization and Montgomery multiplication are
used. We say that a value $x = \sum_{i=0}^{k-1} x_i 2^{si}$ is
\emph{signed-normalized} (respectively \emph{unsigned-normalized}) if
limbs $x_0$ to $x_{k-2}$, i.e. all limbs except the most significant one
$x_{k-1}$, are in the signed $s$-bit range $[-2^{s-1}, 2^{s-1} - 1]$
(respectively, the unsigned $s$-bit range $[0, 2^s - 1]$). The
\verb+normalize()+ function transforms an input integer into another
representation of the same integer value, but normalized. The shown
pseudocode supports both signed and unsigned normalization, depending on
the value of the \verb+signed+ parameter, but any actual implementation
might not necessarily need both. The size-optimized implementations of
ECDSA signature verification shown in sections~\ref{sec:bytecode} use
only unsigned normalization; a variant described in
section~\ref{sec:range:variants} uses both.

Normalization does not modify the integer value; it efficiently contains
limb growth except for the top limb $x_{k-1}$, which can be said to
receive the excess bits. Normalization does nothing about overall value
growth, for which the main tool is Montgomery multiplication, which also
reduces the range of the output, as explained in
section~\ref{sec:montymul:repr}. Section~\ref{sec:range} includes more
details on limb and value growth.

Normalization, as presented in \verb+normalize()+, uses a sequence of
successive operations with dependencies (the value \verb+cc+ is obtained
from the previous addition and needed as input for the next). These
dependencies could be removed by delaying the addition of the upper bits
from the previous limbs (\verb+cc+) to after the truncation of \verb+z+
to $s$ bits, rather than before. In such a case, the normalization is
only partial (output limbs may slightly exceed their nominal range) but
partial normalization could be enough to contain limb growth, and the
removal of dependencies can only help CPUs with out-of-order execution
and SIMD implementations. Full normalization will still be needed for
canonicalization purposes, typically right before encoding output
values.

\paragraph{Pre-shifting.}
The \verb+longmul+ operations return the product over a pair of machine
words; that value must then be split into two limb-sized parts. When the
representation is non-redundant and the radix matches the word size, the
two words returned by \verb+longmul_u+ are these low and high parts, and
the split is implicit and free. This is no longer the case with a
redundant representation. Given a product $c$ which is obtained as two
$w$-bit words $c_0$ and $c_1$, we have $c = c_0 + 2^w c_1$, and we want
the low and high radix-$2^s$ halves $c_L$ and $c_H$ such that $c_L \in
[0, 2^s - 1]$ and $c = c_L + 2^s c_H$. The low half $c_L$ can be
obtained from $c_0$ by masking out the top $w-s$ bits, which can be done
with a bitwise \textsc{and} operation (which is inexpensive); however,
$c_H$ requires combining bits from both $c_0$ and $c_1$. Some CPUs have
dedicated instructions to extract a word value straddling over two
registers (e.g. \verb+shld+ on x86, or \verb+extr+ on Armv8-A), but some
others do not, and computing $c_H$ then needs two shifts and a bitwise
\textsc{or} (RISC-V, as usual, is in that situation). In order to make
this split easier and faster, we use \emph{pre-shifting}:
\begin{itemize}

    \item Operands for a multiplication are left-shifted by some bits
    prior to the multiplication: the limbs of the first operand
    (\verb+a[i]+) are left-shifted by \verb+sh1+ bits (line 49) and
    those of the second operand (\verb+b[i]+) are left-shifted by
    \verb+sh2+ bits (line 46), for two fixed constants \verb+sh1+
    and \verb+sh2+ whose sum is equal to $w-s$. The constants must be
    chosen so that these shifts do not overflow the represented limb
    values, which depends on limb growth.

    \item The net effect of pre-shifting is that the \verb+longmul_s+
    operation returns not the formal product $c$, but $2^{w-s} c$. In
    particular, this implies that the high output word $c_1$ now
    contains exactly $c_H$, so that no complicated word extraction is
    necessary. As for the low word $c_L$, it can be obtained with a
    single right-shift of $c_0$ (the \verb+rsh_u+ operation on line 22).

    \item For the addition of the modulus $m$ multiplied by the
    Montgomery factor, all the pre-shifting is done in the factor,
    while the modulus is unshifted.

\end{itemize}
In total, pre-shifting requires $2k$ extra left shifts for the operands,
but each of the $2k^2$ splits of \verb+longmul_s+ outputs only needs a
single right-shift operation.

\paragraph{Signed Montgomery factor.}
A not immediately obvious feature of the Montgomery multiplication shown
in listing~\ref{alg:mmul-redundant} is that the Montgomery factor
\verb+f+ is signed. Formally, the overall Montgomery factor $f$ (in the
computation of $(ab + fm)/R$) is obtained limb by limb, and each limb
$f_i$ is, in the classic description, a non-negative value (it is in
$[0, 2^s-1]$ when limbs are in radix $2^s$). In the redundant
representation:
\begin{itemize}

    \item The factor $f_i$ (in variable \verb+f+ in the code, line 51)
    is obtained by multiplying the low limb \verb+t[0]+ with the constant
    \verb+m0i_sh+ which is equal to $2^{w-s} (\umod{-1/m}{2^s})$ (that is,
    the Montgomery coefficient $g$, left-shifted by $w-s$ bits). A wrapping
    multiplication (modulo $2^w$) is used. Thus, the output \verb+f+
    really contains an $s$-bit value, but stored in a $w$-bit word and
    left-shifted by $w-s$ bits.

    \item The addition of the modulus multiple is done with the
    sub-function \verb+rr_add_mul()+ (the same function as for adding
    the partial product $a_i b$), which uses \verb+longmul_s+ (line 21).
    This interprets the $w$-bit word \verb+f+ as a signed value, in the
    $[-2^{w-1}, 2^{w-1} - 1]$ range. This value is equal to the formal
    $f_i$ multiplied by $2^{w-s}$, therefore $f_i$ is really in the
    signed $[-2^{s-1}, 2^{s-1} - 1]$ range.

\end{itemize}
The fact that the $f_i$ values are signed is crucial for the range
analysis. It means that the overall $f$ factor cannot be (much) larger
than $R/2$ in absolute value, which makes Montgomery multiplication
output values smaller, and helps with value growth. For instance, if
$m$ is a 256-bit integer, $w = 64$, $k = 5$, and $s = 54$ (values are
over five limbs in radix $2^{54}$), then $R = 2^{ks} = 2^{270}$; if
the input operands $a$ and $b$ are at most $30m$ in absolute value,
then $(ab + fm)/R$ is less than about $0.55m$ (in absolute value),
which allows for up to 54 successive limb-wise additions before value
growth reaches $30m$ again. With unsigned $f_i$ values, Montgomery
multiplication outputs could reach up to about $1.055m$ (with inputs
at most $30m$) and only 28 unreduced additions would be allowed until
a $30m$ value growth is again obtained.

\paragraph{Post-multiplication normalization.}
Montgomery multiplication, as presented in
listing~\ref{alg:mmul-redundant}, uses the extra room in each limb to
accumulate partial results without carry propagation. However, the
resulting values can be quite large, and a normalization step is
normally required after the core Montgomery loop. In
\verb+mod_monymul()+, the \verb+normalize()+ function is called with a
\verb+False+ second parameter, to perform unsigned normalization. A
variant would be to use signed normalization. Signed normalization
yields smaller limbs (in absolute value), and thus allows more
downstream unreduced additions and subtractions to be performed (with
regard to limb growth and potential overflows), until the next
multiplication. However, signed normalization is slightly more expensive
(there is an extra subtraction per limb, line 37), and unsigned
normalization is usually also needed (for integer canonicalization
purposes), so that size-optimized implementations will usually prefer to
use only unsigned normalization throughout.

% ----------------------------------------------------------------------

\newpage
\section{Human vs AI}\label{sec:showdown}

In mid-March 2026, over the course of about 8 days, a human developer
(myself) battled against an AI (Claude Code) in a contest for producing
the most size-optimized implementation of ECDSA signature verification.
On my part, one of the main points of the exercise was to find out
whether the redundant representation with signed Montgomery factors was
novel or not, the assumption being that the AI has extensive knowledge
of everything that has ever been published on the Internet. Claude
access and operation were provided by Steve Weis, who thought that the
experiment was fun and had some spare compute tokens.

\subsection{ECDSA}

We consider ECDSA signature verification only, not key pair nor
signature generation. In particular, all values are public, and there is
no applicable notion of a side-channel information leak. A useful mental
model for a size-optimized ECDSA signature verification code is a system
that needs to verify a signature on its own reflashable firmware at
boot-time; the verification code must then fit in a dedicated ROM block,
hidden inside the CPU, where room is scarce\footnote{In that specific
case, there are in fact some alternatives, e.g. the ROM could contain a
hash function implementation, but the signature verification code could
be loaded dynamically from Flash; it would suffice to verify the hash of
that code against a hash value hardcoded into the ROM space.}.

We use NIST curve P-256. The point coordinates are expressed in the
field of integers modulo $p = 2^{256} - 2^{224} + 2^{192} + 2^{96} - 1$,
a 256-bit prime. The curve order is a prime integer $n$ which is close
to $p$ (but slightly lower); we call \emph{scalars} the elements of the
field of integers modulo $n$. Each curve point (except the special
point-at-infinity) has two coordinates $(x,y)$, which are integers
modulo $p$; they fulfill the curve equation $y^2 = x^3 - 3x + B$ for a
given curve constant $B$. A conventional, fixed generator point $G$ is
given. The private key is a non-zero scalar $d$, and the corresponding
public key is the point $Q = dG$. A public key point is commonly encoded
with the ``uncompressed'' format over 65 bytes: one byte of value
\verb+0x04+, followed by the $x$ and the $y$ coordinates, over exactly
32 bytes each: the encoding of a coordinate is the unsigned big-endian
encoding of the unique integer in $[0, p-1]$ corresponding to the
coordinate value modulo $p$.

A signature is a pair of scalars $(r,s)$. In this paper, we use the
so-called IEEE p1363 format, in which a signature is expressed over
exactly 64 bytes: 32 bytes for each of $r$ and $s$, in that order. Each
scalar is encoded as the unsigned big-endian representation of the
unique integer in $[0, n-1]$ which matches the scalar value.

The signed message $m$ is first hashed with some hash function; only the
hash output is used for signature generation or verification. For curve
P-256, the hash output must have size at least 256 bits (32 bytes), and
only the first 256 bits are used; the SHA-256 hash function is commonly
used for ECDSA with P-256. Signature verification is described in FIPS
186-5 (\cite{Fips186}, section 6.4.2), and entails the following:
\begin{enumerate}

    \item Decode the public point $Q$, verifying that:
    \begin{itemize}
        \item the encoded point size is exactly 65 bytes;
        \item the first byte has value \verb+0x04+;
        \item each provided $x$ and $y$ coordinate is the unsigned
        big-endian encoding of an integer in the $[0, p-1]$ range;
        \item the two coordinates $(x,y)$ fulfill the curve equation
        $y^2 = x^3 - 3x + B$.
    \end{itemize}

    \item Check that the provided hash value has size at least 32 bytes,
    and interpret the first 32 bytes as an integer $e \in [0, 2^{256}-1]$
    (using the unsigned big-endian convention). In the defined contest
    API, we will also require the verifier to enforce a \emph{maximum}
    hash value size of 64 bytes (if it is larger than that, then it is
    probably not a hash value, which means that the calling application
    is using the code wrongly and unsafely).

    \item Check that the signature size is exactly 64 bytes, the two
    halves of which encoding integers $r$ and $s$ which are in the
    $[1, n-1]$ range (take care that this involves checking that the
    integers are lower than $n$, but also that they are both non-zero).

    \item Compute $u = \umod{e/s}{n}$ and $v = \umod{r/s}{n}$ ($s$ is
    invertible modulo $n$ because $0 < s < n$ and $n$ is prime).

    \item Compute the point $W = uG + vQ$. If that point turns out to be
    the point-at-infinity, then the verification stops and returns a
    failure.

    \item Get the $x$ coordinate of point $W$ as a canonical integer $x_W$
    in the $[0, p-1]$ range.

    \item The signature is valid if and only if $\eqvm{x_W}{r}{n}$.

\end{enumerate}
A tricky point is that $x_W$ is obtained as an integer modulo $p$, but
is then reinterpreted as an integer modulo $n$. This is why it is
important to take care that it has been properly canonicalized in the
$[0,p-1]$ range. Moreover, $n < p$, which means that interpretation of
$x_W$ modulo $n$ entails an extra modular reduction (but it can be
performed with a single conditional subtraction of $n$, since $n$ is
very close to $p$). Similarly, the hashed message $e$ is turned into an
integer modulo $n$ but it could be greater than $n$: the implementation
must tolerate an out-of-range value, or first reduce $e$ modulo $n$
(again with a conditional subtraction).

\subsection{The Event}

The formal goal for the implementations is a function that verifies an
ECDSA signature over curve P-256; inputs are the encoded public key
($Q$, over 65 bytes), the signature (64 bytes), and the hashed message
(hash value, 32 to 64 bytes). I defined a simple one-function API in
C, shown in listing~\ref{alg:ecdsa-API}; it includes comments which
define precisely the expected semantics. Notably, the input parameter
sizes are provided explicitly, and must be checked by the function;
also, the hash value length must be verified to be at least 32 bytes,
but also at most 64 bytes. This file was provided to the AI.

\begin{figure}
\begin{forbidpagebreak}
\lstinputlisting[caption={\label{alg:ecdsa-API}API for ECDSA signature verification (over NIST curve P-256)},style=zzc]{ecdsa-p256-verify.h}
\end{forbidpagebreak}
\end{figure}

The implementation target was a 64-bit x86 system (which matches the
virtual machine to which Claude was given access, to run test versions),
assuming a ``bare metal'' model (no standard library available, so
everything had the be done locally, even moving memory blocks).
Compatibility with the standard 64-bit x86 System V ABI (used in Linux
systems under the architecture name ``x86\_64'', though ``amd64'' is
also encountered) was to be respected (the AI turned out to have a bit
of trouble with that rule, it kept trying to use extension opcodes such
as \verb+mulx+ and had to be reined in explicitly).

ECDSA signature verification on curve P-256 is not a novel thing; there
are plenty of implementations around, so that we could assume that
Claude would be able to make running code, if only by reproducing an
existing implementation. Moreover, size-optimization on that particular
platform is known to have lots of ``local tricks'' by leveraging the
complicated instruction encoding, due to the very long history of that
specific instruction set (e.g. even if instructions can work with any of
the registers, the ``historical'' registers yield smaller encodings); we
thus expected a pattern-matching engine like an LLM-based AI to be quite
good at that task. To ensure correction of the result, Steve provided
the AI with a fresh copy of relevant test vectors from project
Wycheproof\cite{Wycheproof}, which were designed to exercise many edge
cases.

When the competition started, I already had a C implementation using
54-bit limbs in a redundant Montgomery representation, which I had
manually translated to assembly, for a code size of about 2080 bytes.
The only non-trivial (but still quite classic) size optimization was the
use of a custom 16-bit instruction format to represent over 86 bytes the
43 operations involved in the complete point addition formulas from
Renes, Costello and Batina\cite{RenCosBat2016}: these are additions,
subtractions and multiplications over a handful of local variables,
which are readily mapped to three-operand instructions (4 bits for each
of two sources and one destination, and 4 bits for the operation
itself). The initial productions from Claude used classic Montgomery
representation (with four 64-bit limbs) and started with non-complete
formulas over Jacobian coordinates. The AI quickly found and adopted its
own 16-bit custom format; after being told by Steve to better look at
existing formulas, it remembered that the Renes-Costello-Batina formulas
exist. More grinding at local optimizations led it to about 1450 bytes.

Meanwhile, my own local optimization efforts had reduced my code to
about 1640 bytes, which was not good enough. However, it allowed me to
gain a finer understanding of what operations were actually needed for
ECDSA signature verification; I knew from the start that a \emph{really}
reduced implementation would have to use a comprehensive custom
instruction set, running in what amounts to an emulated CPU, with very
small instructions, preferably over one byte each (thereafter called
``bytecode''). Designing that instruction set required first knowing
what instructions would be useful for the task at hand. At that point, I
knew enough to attempt it, so I did, and got my code down to 1154 bytes.
Steve conveyed that result to Claude, with some appropriate orders and
taunting to make it grind futher. The AI switched to a non-Montgomery
representation with 32-bit limbs, leveraging the special format of the
involved moduli: for curve P-256, when using 32-bit limbs, then both the
coordinate field modulus and the scalar field modulus (i.e. the curve
order) are such that the most significant limb is equal to $2^{32}-1$,
which allows for a specialized reduction routine.

Claude and myself worked concurrently, finding and applying local
optimizations. I was first to get my code below the symbolic threshold
of 1000 bytes. After a while, the AI got to two implementations at 890
and 908 bytes (the latter being noticeably faster) but was getting
stuck; Steve tried to nudge it toward trying some other representations,
in particular using five 54-bit limbs, but it got nowhere (which was
indirectly a confirmation that the kind of redundant Montgomery
representation that is described in this paper was not in Claude's
training set for ECDSA). At that time, my own code was down to 929 bytes
(I later got it down to 875 bytes, and an alternate version at 848
bytes, but it was in the post-competition cleanup). Since the
optimization efforts were slowing down, I decided to end it with a last
parting shot: a stupidly slow alternative, which does not even use
Montgomery representation at all, but allows to reduce the code size to
784 bytes (which I quickly shrunk to 766, then 732 bytes). The robot was
stumped, and Steve conceded defeat in its name.

\subsection{Why Did Claude Lose?}

The AI can only use pre-existing knowledge (what is on the Internet,
basically) and cannot exploit any new trick that has not been published
beforehand, and is not easily inferrable through the kind of pattern
matching that is at the core of how an LLM works. It is understandable,
for instance, that it could not leverage the notion of a redundant
representation with signed Montgomery factors. However, my final (and
slow) code uses none of these mathematics; it is in fact very basic, and
leverages only well-known methods that should have been accessible to
Claude.

The fact that Claude did not find it is probably an artifact of how it
works: mostly by ``grinding'' incremental optimizations. It quite early
chose a 16-bit instruction format for translating curve point addition
formulas (this is natural, everybody does it --- I did too, in my first
implementations), and then tried to augment it with extra instructions,
but remained stuck on the 16-bit size. By constrast, my final code uses
8-bit instructions, which allow for expressing more behaviour in fewer
bytes. Moving from 16-bit to 8-bit instructions is a structural jump
which at first would make the code quite larger, even though it
ultimately allows for a more compact solution, after sufficient
optimization rounds. This is a kind of ``stepping back''
meta-optimization that humans can do but is apparently harder for AIs;
the way it works for me is that when I am thinking about a particular
optimization problem, I am \emph{also} doing some meta-analysis in the
back of my mind, always trying to fit the task in a larger context and
considering entirely different strategies. This background thinking can
keep on going while I perform other mechanical tasks that do not require
my entire focus\footnote{As it were, I found the main optimization of
that implementation, i.e. removing multiplications entirely, while
folding laundry.}. It seems that the current AIs do not have such
background thinking abilities.

The ``stupid'' code is available in a separate
repository\footnote{\url{https://github.com/pornin/codegolf-ecdsa}};
since it is very expensive to use (runtime cost is around 255 million
cycles!), we will not discuss it further in this paper. In all of the
following, we will focus on the ``non-stupid'' code, which in particular
leverages the improved Montgomery multiplication with a redundant
representation, as explained in section~\ref{sec:montymul:redundant}.

\subsection{Benchmarks and Domination}

Near the start of the battle, I had provided not only the API, but also
a loose notion that the code should not be so slow that it was unusable,
with a fuzzy limit at one billion clock cycles; I also gave a few
figures measured on my own code. Steve conveyed these values to Claude,
which took it as a hint that speed was very important, and it kept
expressing its optimizations with sentences such as ``Thomas is
DOMINATED!'' with uppercase letters, a grim reminder that, as a large
language model, this AI was trained, alas, on the Internet\footnote{The
``domination'' terminology is from the academic field of multi-objective
optimization and relates to Pareto frontiers; but the uppercase taunting
is from less well-behaved areas.}.

Notwithstanding, it was getting some results which looked slightly too
good. After the event, I tried Claude's code on my own system, and found
cost figures roughly $1.5$ times higher than what the AI was reporting.
Looking at how it was measuring its own performance, it turned out that
Claude relied on the ``time-stamp counter'' (\verb+rdtsc+) which is
provided in x86 CPUs, without any regard for frequency scaling. In older
times (say, about 20 years ago), the time-stamp counter was faithfully
counting actual CPU cycles, and \verb+rdtsc+ was a good way to measure
cost of cryptographic primitives. However, CPUs have since deployed
automatic frequency changes to adapt to load but also available voltage
and current temperature; this goes both ways, the CPU being able to run
at a \emph{higher} frequency than its nominal speed, though not
continuously, because that would induce overheating (Intel calls this
feature ``TurboBoost''). The time-stamp counter now runs at a fixed,
nominal frequency matching the CPU baseline but not changing when the
CPU actually runs either slower or faster. The bottom line is that for
every clock cycle that Claude was measuring with \verb+rdtsc+, the CPU
had actually executed $1.5$ cycles worth of instructions, making it
underestimate its own costs by that factor $1.5$. On my own test system,
I use \verb+rdpmc+ to access performance counters, which return the true
cycle count (but require the access to have been allowed by the system
administrator); I also fully disabled TurboBoost\footnote{I disabled
TurboBoost primarily to keep the CPU cooler, which implies that the CPU
fan is running slower and more silently; I hate fan noise. But a
side-effect is that it makes my benchmarks much more stable and
reliable.}.

So much for domination. The robot kept boasting about its superior speed,
though it was, in hindsight, an illusion. It never really obtained better
performance.

\subsection{The Case of Test Vector 481}

During post-event testing of Claude's output, a bigger and somewhat more
concerning surprise was that when I tried it in my own test framework,
it failed on one test vector. This was unexpected because the AI was
working on the Wycheproof test vectors, which my test framework also
uses, and the failed test was on one of these vectors. Was the AI really
running its own tests? The explanation turned out to involve some
historical quirks of ECDSA.

ECDSA, like its predecessor non-elliptic-curve DSA, represents
signatures as a pair of scalars $(r,s)$. The ECDSA standards (FIPS
186-5, ANSI X9:62) do \emph{not} define a standard serialization of such
pairs into bytes, so that several different, incompatible encodings have
emerged. The simplest one is the IEEE p1363 format, which we use here in
our ECDSA implementations; a widespread alternative is to store $r$ and
$s$ into an ASN.1 \textsc{sequence} of two \textsc{integer} values,
which is then DER-encoded into bytes. The ASN.1-based format is the one
used in, for instance, X.509 certificates. ASN.1 and p1363 are the two
main conventions in widespread usage (there are a few others but they
are not as widely used). Correspondingly, project Wycheproof contains
two distinct lists of test vectors, for the p1363 format (file
\texttt{\small ecdsa\_secp256r1\_sha256\_p1363\_test.json}) and for the
ASN.1 format (file \texttt{\small ecdsa\_secp256r1\_sha256\_test.json}).

Since we were using the p1363 format, Steve, quite logically, provided
the AI with a copy of the test vectors in p1363 format. On my side,
being somewhat of a perfectionist\footnote{``Perfectionist'' is a nice
terminology to designate a raving maniac.}, I used \emph{both}: some of
the ASN.1 test vectors specifically exercise DER decoding issues, but
all others use properly ASN.1-encoded signatures, and I could
automatically convert these into the p1363 format, thus making a bunch
of additional test vectors. One of these --- test vector 481 --- was part
of a few test vectors added by David Benjamin in January 2026, designed
specifically to detect a potential issue in an implementation trick, but
David only provided the ASN.1 version, not the p1363 version. Claude
produced code that used this trick, and had the bug, but since it was
working only with the Wycheproof p1363 test vectors, it did not notice
it. My own test framework turned out to be more comprehensive (and I did
not have that issue to begin with).

The bug is related to the final steps of the verification, when the $x$
coordinate of $W$ is extracted as an integer $x_W$, which is then
reinterpreted as an integer modulo $n$. Internally, $W$ uses some kind
of fractional representation, typically projective coordinates, and
reconstructing $x_W$ requires computing a modular division (modulo $p$):
$\eqvm{x_W}{X/Z}{p}$ for two integers $X$ and $Z$. Inversion of $Z$ is
somewhat expensive, both in computation time and in code size. To avoid
these costs, the trick is to assert that an integer $r\in [0, n-1]$ may
correspond to only two possible source integers $x_W\in [0, p-1]$ when
they are reduced modulo $n$; these are $r$ and $r + n$. We can thus
backport the final two steps of ECDSA signature verification into
integers modulo $p$, and implement them by checking whether
$\eqvm{X/Z}{r}{p}$ or $\eqvm{X/Z}{r + n}{p}$, which can be further
simplified into testing that $\eqvm{(X - Zr)(X - Z(r + n))}{0}{p}$.
Hurray, no more division!

This is a nice trick, with the slight drawback that it is, in fact,
wrong, at least in the way it is presented above. A target integer $r\in
[0, n-1]$ can be obtained from \emph{up to two} possible integers in
$[0, p-1]$, but usually only one is possible. If it so happens that $r <
p - n$, then indeed $r + n < p$ and if $x_W$, as an integer, is equal to
either $r$ or $r + n$, then the signature will be valid. However, most
$r$ values are greater than $p - n$, and then $r + n \geq p$ and is no
longer a valid target value for $x_W$. The optimization is applicable
only if $r + n < p$, but some implementations omit that test and
\emph{always} accept $r + n$ as a valid target (modulo $p$), or accept
it if $r + n$ is greater than $p$ but ``not too big'' (e.g. it fits in
256 bits).

Test vector 481 exercises this potential issue: it uses a value $r$
which is very slightly above $p - n$, and a verification point $W$ whose
affine $x$ coordinate is $x_W = r + n - p$ (when canonicalized into $[0,
p-1]$). Strict application of the ECDSA rules leads to the conclusion
that the signature is invalid, but Claude's implementation accepts it as
valid, which probably means that it is using the trick described above
but lacks the required test that $r < p - n$.

Formally speaking, the AI's productions were never fully correct, and
would have to be disqualified from the event. In a less rigid
apprehension, we may consider that Claude leveraged an existing
technique which was also wrong; the fact that a corresponding Wycheproof
test vector exists at all shows that the bug is also encountered in the
wild, and the AI lifted it, not knowing that it is incorrect (or missing
that the range test on $r$ was necessary).

\subsection{Analysis of AI Usability}

The ``artificial intelligences'' based on large language models, which
are highly fashionable right now (however quaint it might seem to future
readers), have a large number of objectionable externalities, such as
the stupendous energy and water usage of large data centres; the
disregard for any intellectual property or even courteous attribution
for the sources that the AI reproduces; the ruthless and clumsy
large-scale scraping of training data that is often indiscernable from
downright denial-of-service attacks; etc. Additionally, programming is
both a technique and an art, and exercises such as extreme code size
optimization (colloquially known as ``codegolf'') are obviously more on
the artistic than technical side; artistically, AI-generated code is
some abhorrent slop, a repellent effusion of ugliness, in the words of
Hayao Miyazaki it is ``an insult to life itself''. One can guess from
this paragraph that I am not overly fond of the whole AI endeavour. Yet
I think that it is possible to examine the technical usability and
utility of the tool in a specific task (here, implementing cryptographic
primitives) dispassionately, independently of these considerations. We
will thus try in this subsection to focus on the technical aspects of AI
coding for cryptographic primitives.

On the plus side, Claude showed a strong ability at finding and applying
local optimization tricks, wading with ease into the corners and nooks
of the arcane encoding of the x86 instruction set. The ability to
maintain a (mostly) functioning ECDSA verifier while renaming registers
and moving functions around is impressive. We initially conjectured that
the AI would be somewhat good at it, and it exceeded my expectations. It
is also quite relentless in the ``grinding'' tasks.

On the minus side, it did make mistakes, as explained above: it
unquestionably used an outdated benchmarking method, and it applied an
incorrect optimization trick. A human developer might make the same
mistakes, so this is not a sufficient reason to dismiss code-generating
AIs entirely. A more interesting question is then: assuming that your
robot did make some mistakes, what do you do? How should that be fixed?
A salient point of the generated assembly code is that it is not exactly
readable. The AI was instructed to add explanatory comments, and they
can be used to somewhat infer how the computation happens in broad
terms, but this is not sufficient to allow manual modifications, e.g. to
fix the issue with test vector 481. Internally, Claude produced and
modified its assembly code through custom code generators in Python,
which adds an extra layer of unintelligibility. The only reasonable path
to fixing the bug is to provide the test vector to the AI and ask it to
fix its code. This raises three noteworthy points:
\begin{itemize}

    \item If you want to modify such AI-optimized code then you need the
    AI again. This has become a dependency. You cannot readily remove
    the AI from your maintenance strategy, even if, for instance, the AI
    company decides to raise usage price by a factor of 10, or in the
    event that the AI company, being a startup, hits a financial wall
    and is wiped out into bankruptcy then liquidation. Human developers
    writing code which is not maintainable by other humans are a known
    issue; AIs are not immune to that, and cannot be assumed to always
    be around.

    \item AIs cannot learn. If a human developer makes a mistake, you
    can talk to that developer about it, and they will remember the
    mistake and how to avoid it; the same human might even proactively
    inform other humans about it, a process known as ``teaching''. This
    cannot be really done with AIs as currently deployed, because their
    ``mind'' is split into two parts: on the one side, during the
    training, the internal LLM was trained, at stupendous cost, on an
    extremely large set (at first order approximation, the entire
    Internet!) and can potentially ``remember'' all or most of it. On
    the other hand, the AI \emph{instance} that you are using has about
    no memory at all. It can \emph{read} all its training (which is
    ingrained into its emulated neural paths) but it cannot modify it.
    The only way for an LLM to remember anything is to put it in its
    local window, which is continuously fed back into the inference
    engine. This memory is entirely local to your specific instance; if
    you fire up another Claude session it will remember none of it. The
    window is also very small, and this is structural to the way it is
    processed (the window is repeatedly reinjected entirely as inference
    produces output tokens). The end result is that if you want an issue
    and its avoidance to become a somewhat permanent piece of knowledge,
    then you \emph{must} turn it into a test vector that should be fed
    to all other AIs working in the same area.

    \item AIs can flood you with generated comments and documentation,
    but that does not necessarily lead to a better understanding of what
    the AI did. A human developer who found a nifty trick is likely to
    boast about it, make blog posts and videos and presentations at
    hacking conferences, and generally trying to demonstrate all the
    elegance of the method and how cunning and aesthetically pleasing it
    is\footnote{This entire paper clearly shows how much guilty I am of
    this particular trait.}. The human will take great care to be
    understood by at least a few other people. You cannot expect that
    from an AI, which does not have any long-term planning and does not
    follow up on its personal communication goals --- it does not have
    personal goals at all. This makes it hard to ensure preservation of
    knowledge. In other words, AIs cannot learn, but they are also bad
    at teaching anything which is not in their immutable training set.

\end{itemize}

Every technical area has its own specificities; for the quite
specialized task of implementing cryptographic primitives, a well-known
feature is that one cannot use functional tests for security. Unit tests
are good at ensuring that a given piece of code behaves normally when
processing normal data. However, cryptography postulates a sentient
attacker who will purposefully try to exercise edge cases, especially
some cases hiding in dark mathematical corners that are outside of reach
of automated strategies such as fuzzing. The only known method to make a
correct implementation is a combination of ``being very careful'' (the
traditional way), and formally proving correctness in a way which is
independently verifiable, most preferably with automated verification.

What the discussion above shows is that AIs are good at producing code
which successfully passes a given set of unit tests. However, unit tests
never offer a complete validation; there can always be some issues that
will not be detected by the tests, even for very good and comprehensive
sets such as Wycheproof. The fast exploration of numerous variants,
which is the very thing that makes AI-generated code technically
desirable, also makes them good at finding new kinds of bugs that bypass
even comprehensive testing strategies. What this means is that when
generating code in an adversarial context (which is the case for
cryptography, by definition), an AI such as Claude should be used only
if it can produce not only the code, but also a verifiable proof that
the code is correct, e.g. using tools such as Lean\cite{Lean}.

Arguably, the same could be said of humans writing cryptographic
implementations. The existence of subtly incorrect code is not a new
problem; AIs only make it more apparent because they are faster at
producing code, often operated without careful human review (AIs are
being deployed \emph{because} they are much faster, and careful
reviewing is the first casualty in a search for speed), and not much
amenable to personal improvements through local learning and teaching.
Using an AI to generate a cryptographic implementation really augments
the importance of having the AI also provide sufficiently detailed
descriptions and proofs of correctness to enable independent validation.

Unfortunately, there is no proving method that is really complete (there
is always some unproven interface between between the physical world and
the world of abstract proofs); however, we can hope for proving
correctness of some complicated parts in order to help a manual review
by a human. Section~\ref{sec:bytecode} describes my own implementation
in details, and in particular how it is structured around a custom
bytecode which is leveraged in section~\ref{sec:range} to perform a
comprehensive and automated range analysis that demonstrates that the
operations on big integers faithfully follow the mathematical
description from section~\ref{sec:montymul:redundant}.

% ----------------------------------------------------------------------

\newpage
\section{Size-optimized ECDSA Verification}\label{sec:bytecode}

\subsection{Implementations and Performance}

I made several implementations, for various architectures:
\begin{itemize}

    \item A pure C implementation is provided; it is portable (no
    unaligned accesses, no reliance on any particular endianness)
    and comes in two flavours, which we designate under the names
    \verb+gen64+ (values over $k = 5$ limbs in radix $2^{54}$)
    and \verb+gen32+ ($k = 12$ limbs in radix $2^{22}$). \verb+gen64+
    is available when the platform supports $64\times 64\rightarrow 128$
    multiplications; \verb+gen32+ uses only $32\times 32\rightarrow 64$
    multiplications.

    \item Optimized assembly implementations for 64-bit x86
    (\verb+amd64+), 64-bit Arm (\verb+arm64+) and 64-bit RISC-V
    (\verb+rv64+). All three use the same 5-limb format as \verb+gen64+,
    with the exact same sequence of operations.

    \item An alternate implementation for 64-bit x86 (\verb+amd64alt+)
    which uses twelve limbs and radix $2^{22}$; it uses \emph{almost}
    the same operations as \verb+gen32+; the exact details are described
    in section~\ref{sec:range:variants}.

\end{itemize}

For size and speed performance, we use three test platforms:
\begin{itemize}

    \item 64-bit x86: Intel i5-8259U at 2.3~GHz (Coffee Lake core);
    TurboBoost is disabled.

    \item 64-bit Arm: Broadcom BCM2712 at 2.4~GHz (Arm Cortex A76 core);
    this is a Raspberry Pi 5 single-board computer (SBC).

    \item 64-bit RISC-V: StarFive JH7110 at 1.5~GHz (SiFive U74 core);
    this is a StarFive VisionFive2 SBC.

\end{itemize}
All three platforms run an up-to-date Ubuntu 24.04 operating system.
Tests are done with C compilers Clang 18.1.3 and GCC 13.3.0, with
compilation flags \texttt{-Os -fno-builtin
-\nobreak fno-asynchronous-unwind-table}; these flags imply optimization
for code size rather than speed, no generation of calls to external
functions such as \verb+memcpy()+, and no generation of stack unwind
tables (these tables improve debugging but otherwise have no effect on
runtime, and they increase code footprint, so we remove them). The
choice of C compiler and its options only impact the C implementations.

Table~\ref{tab:perf} reports the binary code size (in bytes, as reported
by the \verb+size+ tool on the compiled object file) and the average
runtime cost (in millions of clock cycles, measured with the relevant
in-hardware performance counters).

\begin{table}
\begin{center}
{\renewcommand{\arraystretch}{1.1}%
\setlength{\tabcolsep}{1em}%
\begin{tabular}{|l|r|r|r|r|r|r|}
\hline
\textsf{\textbf{Architecture}}
& \multicolumn{2}{c|}{x86}
& \multicolumn{2}{c|}{Arm}
& \multicolumn{2}{c|}{RISC-V} \\
\cline{2-7}
& \multicolumn{1}{c|}{size} & \multicolumn{1}{c|}{cost}
& \multicolumn{1}{c|}{size} & \multicolumn{1}{c|}{cost}
& \multicolumn{1}{c|}{size} & \multicolumn{1}{c|}{cost} \\
\hline
\texttt{gen64} (Clang) & 3019 &  1.68 & 3054 & 3.40 & 2228 &  4.72 \\
\texttt{gen64} (GCC)   & 2565 &  2.25 & 2764 & 3.38 & 2704 &  5.03 \\
\texttt{gen32} (Clang) & 2887 &  6.90 & 3168 & 6.99 & 2478 & 18.73 \\
\texttt{gen32} (GCC)   & 2486 &  7.29 & 2888 & 7.31 & 2950 & 17.84 \\
Assembly               &  875 &  4.59 & 1136 & 4.01 &  984 &  7.05 \\
Assembly (alt)         &  848 & 13.64 &    - &    - &    - &     - \\
\hline
\end{tabular}}
\end{center}
\caption{\label{tab:perf}Binary size (bytes) and runtime cost (millions
of clock cycles) for our optimized ECDSA signature verification code,
on 64-bit architectures.}
\end{table}

We can make a few remarks on the performance figures reported in
table~\ref{tab:perf}:
\begin{itemize}
    \item In general, GCC seems to be slightly better than Clang at
    optimizing for code size. This is however reversed for RISC-V, which
    would point at GCC being not too good, or making no special effort,
    at using the compressed 16-bit instructions (which is normally done
    by prioritizing some specific registers within the compiler's
    register allocation algorithm).

    \item RISC-V turns out to be good for low footprint code. This is
    probably a combination of the extreme RISC-ness of its instruction
    set, making it very friendly to compilers, combined with the 16-bit
    compressed instructions. By contrast, Armv8-A instructions are
    more powerful but use a strict 32-bit size.

    \item The \verb+gen32+ code is normally a bit larger than the
    \verb+gen64+ code because it uses both signed and unsigned
    normalization. On x86, this is counterbalanced by the fact that some
    opcodes that process 32-bit values benefit from shorter encodings
    than their 64-bit counterparts, which is why \verb+gen32+ is somewhat
    smaller than \verb+gen64+ on x86. The same effect explains why
    \verb+amd64alt+ is 27 bytes smaller than \verb+amd64+. Since no such
    effect was expected for Arm and RISC-V, no ``alternate'' assembly
    implementation was done for these platforms.

    \item In general, the Arm Cortex A76 is a powerful core (quad issue
    out-of-order execution with a deep pipeline), on par with
    the Intel Coffee Lake. This shows here as similar performance between
    the two platforms for the compiled \verb+gen32+ implementation. For
    \verb+gen64+, though, the Arm CPU is substantially slower. This is
    due to a misfeature of that specific core with regard to 64-bit
    multiplications with a 128-bit output: these operations need two
    separate instructions which are \emph{not} pipelined, and do not pair
    with each other. The net result is that this core can perform only
    one such long multiplication every 7 clock cycles. By contrast, the
    x86 CPU can fire one such instruction every cycle. The used RISC-V
    CPU is a weaker core (dual issue in-order execution) but it can
    execute one long multipication every 2 clock cycles, which is why
    it is not that much slower than the Arm for \verb+gen64+.

\end{itemize}

The overall speed is of course poor, by usual standards, and even more
so for the assembly code, which uses a custom virtual machine (described
in the next subsections). Internal operations are expressed as loops
which are not unrolled, so that most limb accesses are turned into array
accesses, implying heavy memory traffic and extra index calculations;
all of that is computationally expensive. Moreover, the used algorithms
are not optimized for speed. We can make a rough estimate of the runtime
cost by counting how many multiplications are needed in either the
coordinate or the scalar fields:
\begin{itemize}

    \item Field squarings use the same routine as field multiplication.

    \item All curve point doublings use the generic point addition
    routine (14 multiplications) instead of an optimized point doubling.

    \item The point addition routine is generic for inputs in projective
    coordinates, and does not leverage the fact that in some of the
    calls, the second operand is really in affine coordinates.

    \item Curve operations use a simple double-and-add algorithm with
    no lookup tables and no NAF or wNAF optimizations, so that each
    verification requires an average of 512 calls to the point addition
    routine.

    \item Three inversions are used, using Fermat's little theorem,
    with a basic square-and-multiply that makes more multiplications than
    would be needed (especially for the scalar field $n$: the exponent
    $n-2$ has high Hamming weight).

\end{itemize}
In total, each verification uses about 8400 field multiplications, which
is more than twice what could be expected from a speed-optimized
implementation. However, an important counter-remark about speed for
signature verification is that for a large input, most of the verification
cost is spent hashing the data, not doing the elliptic curve computations,
so that the cost values shown here would probably not matter much in a
practical setup.

\subsection{Virtual Machine}

The size-optimized implementations in assembly are organized as a
sequence of instructions for a custom virtual machine which just a few
instructions, tailored at implementing the needed steps for that
specific task. Instructions fit on one byte each (except for
\verb+_CALL+, which uses two bytes); the sequence of instructions is
thus called here ``bytecode''. The bytecode is very compact and allows
to pack the needed operations in a small size (217 bytes). To this we
must add some constants (136 bytes for the encoded moduli $p$ and $n$,
the precomputed $2^{10}(\umod{-1/n}{2^{54}})$, and the conventional
generator point coordinates), and the code which emulates the virtual
CPU on the concrete hardware, as well as the data decoding routine to
handle the initial input.

At the source code level, the bytecode is written with symbolic
constants and macros, which abstract away the details of instruction
encoding. In order to avoid conflicts between these macros and the
assembler normal mnemonics (e.g. we do not want our ``add'' instruction
to be confused with the CPU opcode \verb+add+), all the symbolic
constants and macro for the bytecode start with an underscore character.
We reproduce this convention in the text below.

The virtual machine offers 32 variables, numbered from 0 to 31. Each
variable can hold a big integer in redundant representation ($k$ limbs
in radix $2^s$). Variable 0 is called the \emph{accumulator} and is the
destination, and one of the operands, of most computational operations.
Other variables are accessed through symbolic constants (at the source
code level) which abstract away the variable indices; some are initialized
to some specific values in the VM startup code sequence:
\newenvironment{codedesc2}
    {\begin{list}
        {}
        {\setlength{\labelwidth}{3.5em}
         \setlength{\labelsep}{0.5em}
         \setlength{\itemsep}{1ex}
         \setlength{\leftmargin}{4em}
         \setlength{\rightmargin}{0em}
         \let\makelabel=\makecodedesclabel
        }
    }
    {\end{list}}
\begin{codedesc2}

    \item[\_M] Current modulus (read-only, returns $p$ or $n$ in
    unsigned-normalized format; default modulus is $n$).

    \item[\_ONE] Constant integer 1.

    \item[\_ZERO] Constant integer 0.

    \item[\_Gx] Generator point, $x$ coordinate, Montgomery
    representation ($\mont{G_x}$).

    \item[\_Gy] Generator point, $y$ coordinate, Montgomery
    representation ($\mont{G_y}$).

    \item[\_Kr] First signature half ($r$).

    \item[\_Ks] Second signature half ($s$).

    \item[\_Ke] Hash value, as an integer ($e$).

    \item[\_Qx] Public key point, $x$ coordinate ($Q_x$).

    \item[\_Qy] Public key point, $y$ coordinate ($Q_y$).

\end{codedesc2}

All these values are unsigned-normalized. The \verb+_Kr+, \verb+_Ks+,
\verb+_Ke+, \verb+_Qx+ and \verb+_Qy+ are freshly decoded from the input
data, and have not been checked yet, i.e. they can be any integer in the
$[0, 2^{256} - 1]$ range; it is up to the bytecode to verify that they
are in their proper ranges and that for instance the public key point is
a real curve point.

Other symbolic names are defined for most other variables, which are
initialized to zero:
\begin{itemize}

    \item \verb+_Wx+, \verb+_Wy+ and \verb+_Wz+ are meant to be the
    projective coordinates of a work curve point $W$. Similarly,
    \verb+_Hx+, \verb+_Hy+ and \verb+_Hz+ are the projective coordinate
    of a second curve point $H$. $W$ and $H$ will be used as the
    operands of the point addition routine.

    \item \verb+_Bm+ will receive $\mont{B}$ (Montgomery representation
    of the curve constant $B$). It is not initialy set when the VM
    starts; the bytecode recomputes it from the generator coordinates
    (this saves space compared to having $\mont{B}$ as a hardcoded
    constant).

    \item \verb+_T0+ to \verb+_T9+ are temporaries.

\end{itemize}

\newenvironment{codedesc3}
    {\begin{list}
        {}
        {\setlength{\labelwidth}{6em}
         \setlength{\labelsep}{0.5em}
         \setlength{\itemsep}{1ex}
         \setlength{\leftmargin}{6.5em}
         \setlength{\rightmargin}{0em}
         \let\makelabel=\makecodedesclabel
        }
    }
    {\end{list}}

Values can be copied between the accumulator (value 0) and other variables
with the following instructions:
\begin{codedesc3}

    \item[\_LD x] Copy the contents of variable \verb+x+ into the accumulator.

    \item[\_ST x] Copy the contents of the accumulator into variable \verb+x+.

\end{codedesc3}

Arithmetic operations use the accumulator as first (or unique) source
operand, and as destination operand; some of them accept a second source
operand, which is provided as an integer argument (the variable index,
normally one of the defined symbolic constants):
\begin{codedesc3}

    \item[\_ADD x] Add the contents of variable $x$ to the accumulator
    (limb-wise addition).

    \item[\_SUB x] Subtract the contents of variable $x$ from the accumulator
    (limb-wise subtraction).

    \item[\_MUL x] Multiply the accumulator with the contents of variable
    $x$. This is Montgomery multiplication; in other words, it assumes that
    the operands are in Montgomery representation. The currently set modulus
    is used.

    \item[\_MRED] Apply Montgomery reduction on the accumulator. This is
    a source code level alias for \verb+_MUL _ONE+ (Montgomery multiplication
    with the integer 1).

    \item[\_NORM] Unsigned normalization of the accumulator.

    \item[\_SMOD y] Set the modulus to either $n$ and $p$, depending on
    the value of the parameter \verb+y+. This value \verb+y+ is
    \emph{not} a variable index, but a conventional designation of the
    relevant modulus. Correct values of \verb+y+ can vary depending on
    the target architecture; thus, this instruction is normally used
    through the aliases \verb+_MODN+ and \verb+_MODP+, which set the
    modulus to $n$ and $p$, respectively. The current modulus is used in
    \verb+_MUL+ and can be read from variable \verb+_M+.

\end{codedesc3}

A single loop facility is provided. These loops do not nest: only a single
loop may be active at any given time. Each loop performs exactly 256
iterations, and internally maintains a loop counter that goes from $255$
down to $0$.
\begin{codedesc3}

    \item[\_FOR] Start the loop. The counter is reset to $255$, and the
    loop iteration point is set to the instruction that immediately follows
    the \verb+_FOR+ instruction.

    \item[\_NEXT] If the counter is not $0$, then it is decremented, and
    the instruction pointer is set to the current loop iteration point
    (just after the \verb+_FOR+ that started the current loop). If the
    counter is $0$, this instruction does nothing. Note that there could
    be several \verb+_NEXT+ instructions for the same \verb+_FOR+.

    \item[\_FNEXT] If the counter is exactly $255$ then \verb+_FNEXT+ has
    the same effect as \verb+_NEXT+. Otherwise, it does nothing. This
    instruction is used to skip the first iteration of a loop.

\end{codedesc3}

Conditional execution uses three ``skip'' instructions: for each, if the
represented condition is true, then the instruction that immediately
follows is skipped (execution ``jumps over'' that following
instruction); otherwise, the instruction that follows is executed.
\begin{codedesc3}

    \item[\_SKIPBITZ x] Skip the next instruction if bit $c$ of the value
    in variable \verb+x+ is zero. $c$ is the current loop counter value.

    \item[\_SKIPNZ] Skip the next instruction if the accumulator contents
    are non-zero.

    \item[\_SKIPNEG] Skip the next instruction if the accumulator contents
    are negative.

\end{codedesc3}
\verb+_SKIPBITZ+ is the only instruction that really uses the loop
counter $c$; it is of course meant to implement square-and-multiply and
double-and-add algorithms. Since this uses the binary representation
of the limbs in memory (this really tests bit $\umod{c}{s}$ of limb
$\lfloor c/s\rfloor$), this assumes that the test variable \verb+x+ is
in unsigned-normalized format. Similarly, \verb+_SKIPNZ+ really tests
whether all limbs are zero, and \verb+_SKIPNEG+ only looks at the sign
of the most significant limb (limb $k-1$): the accumulator contents
must be unsigned-normalized for these tests to mathematically make sense.

A subroutine facility is provided:
\begin{codedesc3}

    \item[\_CALL label] Call the subroutine that starts at the provided
    assembly-level \verb+label+.

    \item[\_RET] Return from the current subroutine.

\end{codedesc3}
In the instruction encoding, an unsigned 1-byte offset is used, so that
the subroutine must be located \emph{after} the call site in the address
space (this prevents recursion). Called subroutines may themselves call
other subroutines (no nesting limitation). There is no facility to pass
arguments, the caller and the callee must conventionally agree on which
of the VM variables are used for the subroutine parameters.

Finally, two instructions exit the virtual machine unconditionally,
with either a success status (signature is valid) or a failure status
(signature is not valid):
\begin{codedesc3}

    \item[\_OK] Exit the VM with a success.

    \item[\_FAIL] Exit the VM with a failure.

\end{codedesc3}

\subsection{ECDSA Bytecode}

The VM instructions are enough to implement all needed operations for
ECDSA signature verification. When the VM starts, the calling code
has already performed initialization steps, which entail the following:
\begin{itemize}

    \item The source signature was verified to have length exactly
    64 bytes. The two halves $r$ and $s$ have been decoded (as plain
    256-bit integers) into the \verb+_Kr+ and \verb+_Ks+ variables;
    the values are unsigned-normalized and non-negative, but at this
    point they may be zero, and they may be as large as $2^{256}-1$.

    \item The public key was verified to have length exactly 65 bytes,
    and its first byte to have value \verb+0x04+. The remaining 64 bytes
    have been decoded into \verb+_Qx+ and \verb+_Qy+ (values are
    unsigned-normalized and non-negative, but not otherwise checked).

    \item The hash value was verified to have length at least 32 bytes
    and at most 64 bytes, and the first 32 bytes have been decoded into
    \verb+_Ke+ (unsigned-normalized, non-negative).

    \item All other variables have been initialized as per the VM
    semantics (current modulus is $n$; variables \verb+_ZERO+,
    \verb+_ONE+, \verb+_Gx+ and \verb+_Gy+ are set to their respective
    values, all others are initialized at zero).

\end{itemize}

Listing~\ref{alg:check_scalar} shows the \verb+check_scalar+ subroutine,
whose role is to verify that the value currently in the accumulator is
a proper input scalar value, i.e. an integer strictly in the $[1,n-1]$
range. This subroutine assumes that the current modulus is $n$ and that
the value in the accumulator is unsigned-normalized. As this is part of
some assembly source code, we can use comments (introduced by \verb+#+)
and we can also use the semicolon (\verb+;+) to separate instructions
that are specified on the same source line. The first \verb+_SKIPNZ+
will skip the subsequent \verb+_FAIL+ if the value in the accumulator
is non-zero: this is the check that the scalar is not equal to zero.
The next three instructions verify that subtracting the current modulus
yields a negative integer; note the use of unsigned-normalization to
make sure that the top limb reliably indicates the sign of the overall
value. If both \verb+_FAIL+ instructions were skipped, then \verb+_RET+
is reached and the subroutine returns to the caller.

\begin{figure}
\begin{forbidpagebreak}
\lstinputlisting[caption={\label{alg:check_scalar}Bytecode: verification that an input value is a proper non-zero scalar.},style=zzpython]{bytecode-check_scalar.s}
\end{forbidpagebreak}
\end{figure}

Listing~\ref{alg:canonicalize} takes an input value (in the accumulator)
which is assumed to be a Montgomery representation $\mont{x}$ of some
value $x$, and it returns $x$ in canonical representation (value is in
$[0, m-1]$ if the current modulus is $m$, and it is
unsigned-normalized). This subroutine inherently leverages the fact that
we use signed Montgomery factors: since the $(f_i)$ are signed values in
$[-2^{s-1}, 2^{s-1} - 1]$, the absolute value of $f$ is such that:
\begin{equation*}
    |f| \leq 2^{s-1} \sum_{i=0}^{k-1} 2^{si}
\end{equation*}
where the limit is equal to $\gamma R$ for some constant $\gamma \approx
1/2$. The main result is that after Montgomery reduction, the output
value may be negative, but in absolute value it will be lower than the
modulus $m$, and it can be canonicalized into $[0, m-1]$ with a single
conditional addition; furthermore, the relevant test in that case is
whether the value is negative, which is readily performed with
\verb+_SKIPNEG+. In the range analysis (section~\ref{sec:range}) we will
also prove that our canonicalization always works (i.e. always returns a
value in the expected $[0, m-1]$ range).

\begin{figure}
\begin{forbidpagebreak}
\lstinputlisting[caption={\label{alg:canonicalize}Bytecode: Montgomery reduction with canonicalization.},style=zzpython]{bytecode-canonicalize.s}
\end{forbidpagebreak}
\end{figure}

Listing~\ref{alg:invert_mod} shows how modular inversion is computed: this
is a simple use of Fermat's little theorem: the input (accumulator) is
raised to the power $m-2$, for modulus $m$. The exponent is set in variable
\verb+_T1+, by subtracting 2 from the modulus; one optimization here is that
for both used moduli $p$ and $n$, this limb-wise subtraction will not make
the value unnormalized (the lowest limb of both $p$ and $n$ happens to be
larger than 2) so we do not need an extra \verb+_NORM+ on the modulus.

\begin{figure}
\begin{forbidpagebreak}
\lstinputlisting[caption={\label{alg:invert_mod}Bytecode: modular inversion.},style=zzpython]{bytecode-invert_mod.s}
\end{forbidpagebreak}
\end{figure}

An important point here is that since all multiplications are Montgomery
multiplications, this really performs inversion on Montgomery
representations: given $\mont{x}$, the routine returns
$\mont{\umod{1/x}{m}}$. A basic square-and-multiply algorithm would
enumerate all bits of the exponent, and start with a value set to
$\mont{1}$. We do not have the constant $\mont{1}$ in the virtual
machine (it was not part of the constants initially set). Instead, we
start the square-and-multiply with the input $\mont{x}$ itself, and skip
the top bit (bit $255$) of the exponent, under the assumption that this
exponent bit has value 1 (which is the case for both moduli $n$ and
$p$). This is the reason why the \verb+_FNEXT+ instruction exists: it
was defined precisely to support this skip of the first iteration in
computing inversions. Supporting \verb+_FNEXT+ in the emulator requires
less code space than embedding the constant $\mont{1}$.

\begin{figure}
\begin{forbidpagebreak}
\lstinputlisting[caption={\label{alg:bytecode-part1}Bytecode: entry sequence.},style=zzpython]{bytecode-part1.s}
\end{forbidpagebreak}
\end{figure}

With \verb+check_scalar+, \verb+canonicalize+ and \verb+invert_mod+
defined, we can look at the start of the main bytecode sequence, which
is shown in listing~\ref{alg:bytecode-part1}. The current modulus is $n$
(this was set during VM initialization) so all operations here work
on scalars. This sequence first checks that the decoded $r$ and $s$
values are in $[1, n-1]$ by using \verb+check_scalar+ on both. Then $s$
is inverted (modulo $n$) and that inverse is multiplied with $e$ and
$r$, yielding two values $u$ and $v$ which are canonicalized (and stored
into \verb+_Ke+ and \verb+_Ks+, respectively --- \verb+_Kr+ is
unmodified).

As explained previously, \verb+invert_mod+ really works on Montgomery
representations, and there is no apparent conversion to Montgomery
representation here. Indeed, such a conversion would be done with a
\verb+_MUL+ with the constant $\umod{R^2}{n}$ (as explained in
section~\ref{sec:montymul:repr}), but that constant is not provided to
the VM. It so happens that such a conversion is not needed here. An
integer $x$ is also the Montgomery representation of $\umod{x/R}{n}$;
thus, the inputs $r$, $s$ and $e$, stored in \verb+_Kr+, \verb+_Ks+ and
\verb+_Ke+ respectively, are also the values $\mont{r/R}$, $\mont{s/R}$
and $\mont{e/R}$. The call to \verb+invert_mod+ on $\mont{s/R}$ computed
the inverse over Montgomery representations, i.e. $\mont{R/s}$. Then,
the \verb+_MUL+ call with \verb+_Ke+ (which contains $\mont{e/R}$)
yielded $\mont{R/s} * \mont{e/R} = \mont{e/s}$. The \verb+canonicalize+
routine uses Montgomery reduction (\verb+_MRED+) to convert the input
out of Montgomery representation, and thus yields $\umod{e/s}{n}$ which
is exactly the formal value $u$ that is needed for the rest of ECDSA
verification. Similarly, the \verb+_MUL+ instruction between $r$ and
the inverse of $s$ computes $\mont{r/R} * \mont{R/s} = \mont{r/s}$ and
canonicalization produces the correct value $v$. At no point is explicit
conversion to Montgomery representation needed; it is done ``naturally''
as a side-effect of inversion through Fermat's little theorem. We
therefore do not need the constant $\umod{R^2}{n}$.

As will be shown below, a similar optimization applies to curve points,
thanks to the use of projective coordinates. The constant
$\umod{R^2}{p}$ can also be omitted. Not needing these two constants
saves at least 64 bytes in the overall implementation
size\footnote{Amusingly, some of Claude's comments on its own code
listed the space savings of not including the $R^2$ constants as the
main motivation for not using Montgomery multiplication; as it turns
out, they are not needed anyway.}.

At this point, values $u$ and $v$ have been computed and are in
canonical format, which is appropriate for testing individual bits in a
double-and-add algorithm for curve point multiplications. The next part
of the main bytecode is shown in listing~\ref{alg:bytecode-part2}: it
switches the current modulus to $p$, then proceeds to compute the value
$\mont{B}$ (Montgomery representation of the curve constant $B$) by
using the curve equation and the provided generator point coordinates.
The generator coordinates were set in Montgomery representation during
the VM initialization, but we will afterwards need them in
non-Montgomery representation, so they are converted with Montgomery
reduction (\verb+_MRED+).

\begin{figure}
\begin{forbidpagebreak}
\lstinputlisting[caption={\label{alg:bytecode-part2}Bytecode: computations of curve constants.},style=zzpython]{bytecode-part2.s}
\end{forbidpagebreak}
\end{figure}

\begin{figure}
\begin{forbidpagebreak}
\lstinputlisting[caption={\label{alg:check_point}Bytecode: verification of the public key point.},style=zzpython]{bytecode-check_point.s}
\end{forbidpagebreak}
\end{figure}

The public key point $Q$ has not been verified yet: the
two provided coordinates $(Q_x, Q_y)$ must be checked to be in the
proper range $[0, p-1]$, and to fulfill the curve equation $y^2 = x^3 -
3x + B$. This is done by the subroutine \verb+check_point+, which is
shown on listing~\ref{alg:check_point}. Note the use of \verb+_MRED+
(line 35) before comparing a value with zero: since \verb+_SKIPNZ+
actually checks that all limbs are zero, this requires ensuring that the
only possible representation of zero (modulo $p$) is the all-zeros limb
sequence. This is ensured by Montgomery reduction, as a side-effect of
using signed Montgomery factors: the output of Montgomery reduction is
normalized, and while it is not \emph{fully} reduced (some values modulo
$p$ accept several representations), it still is in a range $[-\gamma p,
+\gamma p]$ for some $\gamma < 1$, which implies that only one
representation is possible for the modular zero. We could have used
\verb+_CALL canonicalize+ here, but a \verb+_CALL+ is two bytes, while
\verb+_MRED+ is only one.

With the help of \verb+check_point+, we can proceed with the rest of the
ECDSA verification code, as shown on listing~\ref{alg:bytecode-part3}.
It starts with calling \verb+check_point+ to verify the public key
point $Q$; then it computes the point $W = uG + vQ$ with a basic
double-and-add (combined) algorithm:
\begin{enumerate}

    \item $W$ is initialized to the curve neutral element (the
    point-at-infinity).

    \item For $i = 255$ down to $0$:
    \begin{enumerate}

        \item $W \leftarrow W + W$ (point doubling)

        \item If bit $i$ of $u$ is non-zero, then: $W \leftarrow W + G$

        \item If bit $i$ of $v$ is non-zero, then: $W \leftarrow W + Q$

    \end{enumerate}

\end{enumerate}
The \verb+_FOR+ loop and the \verb+_SKIPBITZ+ instruction are used here
to perform these operations. All curve point additions and doubling use
the \verb+point_add_to_W+ routine, which we'll detail later on. Once $W$
has been computed, its $x$ coordinate ($x_W$) is extracted, which
requires a modular inversion, since projective coordinates are used:
\verb+_Wx+ and \verb+_Wz+ contain, respectively, the Montgomery
representations of values $X$ and $Z$ such that $x_W = X/Z$. A call to
\verb+invert_mod+, then a \verb+_MUL+, and finally canonicalization
yield $x_W$ in the format required for correct ECDSA processing:
non-Montgomery representation, and canonical in the $[0, p-1]$ range.

\begin{figure}
\begin{forbidpagebreak}
\lstinputlisting[caption={\label{alg:bytecode-part3}Bytecode: main loop and exit sequence.},style=zzpython]{bytecode-part3.s}
\end{forbidpagebreak}
\end{figure}

Once $x_W$ is obtained as an integer, the rest of the code switches back
to the modulus $n$, and subtracts $r$ (which is still in \verb+_Kr+).
Testing this value against zero is done with a call to \verb+_MRED+,
which, as previously explained, ensures that only the all-zeros
representation is possible for a modular zero, thus allowing
\verb+_SKIPNZ+ to reliably test the value.

A non-trivial optimization, which is explained on lines 36-39, is that
we do not need to handle specially the case of $W$ being the
point-at-infinity. As per the ECDSA standard, if $W$ is the
point-at-infinity, then the signature is invalid and must be rejected.
This is what happens here: in that situation, $Z = 0$, and inversion
(using FLT) returns zero, so that the recomputed $x_W$ coordinate ends
up being zero (modulo $p$, but also as a canonical integer in $[0,
p-1]$). Since value $r$ was previously verified to be non-zero, this
implies that $x_W$ cannot be equal to $r$ modulo $n$, thus leading to a
\verb+_FAIL+ instruction, which is exactly what we want here.

The final element of the bytecode is the \verb+point_add_to_W+ routine,
which is used both for point doublings, and for adding either $G$ (the
curve generator) or $Q$ (the public key). This routine is shown on
listing~\ref{alg:bytecode-point_add_to_W}. It uses projective
coordinates $(X{:}Y{:}Z)$ to represent the affine point $(X/Z, Y/Z)$.
The point-at-infinity uses $(0{:}Y{:}0)$ for any $Y\neq 0$; for all
non-infinity points, $Z\neq 0$. The code is a straightforward
translation of the Renes-Costello-Batina formulas (\cite{RenCosBat2016},
algorithm 4) which are complete for curve P-256, since that curve has no
point of order 2. The source formulas are expressed as three-operand
operations (two sources, one destination), each being mechanically
converted into three instructions (load first source operand into the
accumulator, do the operation with the second source operand, store back
the result into the destination operand), which would lead to a bytecode
size of 130 bytes (including the final \verb+_RET+). However, in many
places the result of an operation is immediately used as source for
another, which allows us to remove some redundant loads and stores;
\emph{in fine}, the \verb+point_add_to_W+ routine uses 85 bytes of code
space, which is actually \emph{shorter} than the 86 bytes needed for
16-bit instructions: even for this routine where 16-bit opcodes seem
most adequate, the 8-bit opcodes are more size-optimized.

\begin{figure}
\begin{forbidpagebreak}
\lstinputlisting[caption={\label{alg:bytecode-point_add_to_W}Bytecode: curve point addition (projective, complete).},style=zzpython]{bytecode-point_add_to_W.s}
\end{forbidpagebreak}
\end{figure}

One final noteworthy point here is how the $\umod{R^2}{p}$ constant for
conversions to Montgomery representation is not needed. We use points in
projective coordinates $(X{:}Y{:}Z)$, and, for any $\lambda \neq 0$,
$(\lambda X{:}\lambda Y{:}\lambda Z)$ is another valid representation of
the same point. Thus, when we obtain affine, non-Montgomery coordinates
$(x,y)$ (for point $Q$ or for point $G$), then the $(x{:}y{:}1)$
projective coordinates are also $(\mont{x/R}:\mont{y/R}:\mont{1/R})$,
which designate the same point, and are all in Montgomery
representations. In other words, the use of projective coordinates gets
us a conversion to Montgomery representation for free. This is also why
we converted $G_x$ and $G_y$ (the generator coordinates) \emph{out} of
Montgomery representation (in listing~\ref{alg:bytecode-part2}, lines
11-14): we needed to complete the affine coordinates into projective
coordinates with the proper $Z$, but we do not have $\mont{1}$, only $1$
inside the VM.

At the end of the computation, we obtain $\mont{x_W}$ but canonicalization
gets us the plain, non-Montgomery value, as befits ECDSA.

\subsection{Assembly Optimizations}

The virtual CPU must still be implemented, which is a matter of local
assembly optimization. Conceptually, the emulator is just a loop which
reads the next instruction byte, then jumps to the code sequence that
handles specifically that byte. The instruction encoding was designed
to support a two-level approach:
\begin{itemize}

    \item A first dispatch is done on the low three bits of the
    instruction (which uses a lookup table with 8 entries). Instructions
    with a numeric parameter (e.g. \verb+_LD+ or \verb+_MUL+) store that
    parameter in the five high bits of the instruction, so they are
    directly handled that way. The address in RAM of the variable
    pointed to by that parameter is systematically computed into a
    dedicated register prior to the dispatch.

    \item Parameter-less instructions (e.g. \verb+_NORM+ or \verb+_FAIL+)
    all have their three low bits equal to 111 (in binary): the corresponding
    first-level dispatch routine then does a secondary dispatch, based on
    the five higher bits.

\end{itemize}

\paragraph{Generic Approach.}
The main general approach is to define a custom call convention. The
instruction-processing routines called by the virtual CPU dispatch are
functions, but they do not need to follow the conventional ABI with
caller-saved and callee-saved registers: these functions are never
called from the outside. \emph{State invariants} such as proper stack
alignment must be maintained (to support asynchronous code execution,
i.e. ``signals''), but we can define the role and preservation of any
register as we see fit. In particular:
\begin{itemize}

    \item Some dedicated registers should be reserved, if possible, for
    maintaining the instruction pointer, the state pointer (where
    variables are located), the current loop context (loop iteration
    point and counter), and the current modulus. Note that the modulus
    is not only its value as an integer, but also the Montgomery
    coefficient $g = \umod{-1/m}{2^s}$, which we store in pre-shifted
    format after the modulus limbs (this is why the modulus is handled
    as a special variable and has dedicated modulus-switching
    instructions).

    \item Routines may share some code. Even if each instruction routine
    is invoked as a function, it may ``tail-call'' into another routine
    with a jump, or even by ``falling-through'' into that other routine
    if it is located immediately afterwards in the code space. In
    our implementations, for instance, all three ``skip'' instructions
    share the same sequence to enact the skip, if it was decided that it
    must be performed.

    \item The native stack is used to support the \verb+_CALL+ and
    \verb+_RET+ opcodes. This implies a non-standard stack management,
    which has a tendency to make debugging somewhat obscure (we do not
    maintain frame pointers, so the debugger cannot print stack traces).

\end{itemize}

The input values ultimately are five big integers, from three distinct
sources (public key, signature, hash value), that must be decoded into
limb sequences. A dedicated function (in our code, we call it
\verb+decode_int()+) should be used and invoked repeatedly; we also use
it for the hardcoded constants (curve generator coordinates, moduli $p$
and $n$), which allows us to store these constants in a compact 32-byte
format, rather than the 40 bytes for five limbs. For some architectures,
the opcode for a local function call is somewhat large; e.g. on x86 in
64-bit mode, a \verb+call+ to a function by name yields a 5-byte
encoding. In these cases, it is beneficial to first store the address of
the \verb+decode_int()+ function in a register, then do the call with
register dereferencing (2-byte encoding per call on x86).

In our implementations, we define the virtual CPU loop as three macros
(initialization, main loop, exit sequence) so that extra functions may
be defined (with conditional compilation) for unit tests that invoke
bytecode sequences. A comprehensive set of such unit tests, in
particular for Montgomery multiplications, inversions, and curve point
additions, is highly recommended.

\paragraph{x86 Optimizations.}
The x86 architecture has a very long history, starting with the Intel
8086, first released in 1978, and already inheriting some design
elements from the earlier 8080 (this was made to allow easier conversion
of existing 8080 code to the new platform). This implies a very
complicated instruction encoding, with various optional prefixes that
have been added along the years. The main consequence is that not all
registers are equivalent, some yielding substantially shorter instruction
encodings. The main elements to remember are the following:
\begin{itemize}

    \item The ``historical registers'' (\verb+eax+, \verb+ebx+, \verb+ecx+,
    \verb+edx+, \verb+esi+, \verb+edi+, \verb+ebp+, and \verb+esp+)
    are usually better; e.g. \verb+addl %eax, %ebx+ (addition in AT\&T
    syntax) is encoded over two bytes, but \verb+addl %r8d, %ebx+, which
    uses the ``new'' register \verb+r8d+, needs three bytes. Some
    specific operations (in particular \verb+and+ and \verb+test+) have
    a further shortened encoding when the first operand is \verb+eax+
    (or its low byte \verb+al+).

    \item 64-bit operations need an extra prefix, thus 8-bit or 32-bit
    operations should be preferred when possible.

    \item Immediate constants that do not fit in $[-128, 127]$ are
    usually encoded over four bytes, and should be avoided. This also
    applies to offsets in address calculations, and to jumps (jumps
    to nearby locations have a smaller encoding).

    \item 32-bit operations automatically zero-extend the destination
    register to 64 bits. To set a register to zero, use \verb+xorl+
    (32-bit bitwise \textsc{xor}) on the register with itself. This is
    compact (2-byte encoding) and is also semi-officially ``blessed''
    (this sequence is recognized by the CPU instruction decoder, and
    handled in the register renaming unit, for an approximately null
    runtime cost).

    \item ``Old'' instructions have the shortest encodings. This
    includes instructions such as \verb+loop+ (decrement a counter and
    jump if non-zero) and \verb+xlat+ (table lookup with a 1-byte
    index). These instructions are often slow (microcoded execution) and
    have strict constraints on the involved registers, which impacts
    register allocation.

    \item Some instructions perform a memory access \emph{and} an index
    update. These are in particular the \verb+lods+, \verb+stos+ and
    \verb+movs+ families of instructions. Judicious use of these
    instructions can help reduce code footprint, but they have strict
    requirements on register allocation (e.g. \verb+lods+ always uses
    \verb+rsi+ as index, and loads the value into
    \verb+eax+/\verb+rax+).

    \item \verb+push+ and \verb+pop+ are encoded over a single byte each
    when the involved register is an historical one, and move 64 bits at
    a time. To copy a 64-bit value from an historical register to
    another one (e.g. \verb+rcx+ to \verb+rdi+), a
    \verb+push+/\verb+pop+ combination is two bytes, while a 64-bit
    \verb+mov+ is three bytes. Note that the \verb+mov+ is faster, since
    it is usually handled by the register renaming unit, and involves no
    memory traffic (the CPU will normally do the memory write of the
    \verb+push+ asynchronously, and optimize the \verb+pop+, but the
    memory write will still need to happen at some point, which is not
    free).

\end{itemize}

\paragraph{RISC-V Optimizations.}
RISC-V defines core integer instructions (I) and many sets of extra
optional instructions; any hardware implementation may include some of
these extra instruction sets, and formal \emph{profiles} are defined
which list instruction sets whose presence is guaranteed. For our ECDSA
implementation, we use the 64-bit mode with the core integer
instructions (I) as well as the multiplication instructions (M --- only
multiplication instructions, not division), and the compressed
instruction encodings (C). These instruction sets are part of all
currently defined ``application-processor'' profiles.

RISC-V instructions are simple, with a regular encoding; most of the
size optimization efforts are in register allocation to benefit from the
compressed instruction encodings (16-bit instructions instead of
32-bit): while the CPU has 31 general-purpose registers which are mostly
interchangeable, most compressed encodings can be used only with the 8
``popular'' registers (\verb+s0+, \verb+s1+, and \verb+a0+ to
\verb+a5+). To achieve a low code size, the most used registers should
be chosen among these registers.

Another noteworthy point is that the local function call instruction
(\verb+jal+) with the target being designated by a label uses a 32-bit
encoding, while the call-to-register instruction (\verb+jalr+) can use a
compressed encoding. Just like on x86, we leverage the shorter
\verb+jalr+ by first obtaining the address of the \verb+decode_int()+
function into a register, then using it repeatedly with \verb+jalr+.
Similarly, getting the address of a function or a static variable into a
register uses the \verb+lla+ sequence, which consists of two 32-bit
instructions; we endeavour to have as few such sequences as possible,
hence we ensure that \verb+decode_int()+, the block of constants, and
the bytecode entry point are all consecutive in the address space, so
that a single \verb+lla+ may be shared for all three elements.

\paragraph{Armv8-A Optimizations.}
Armv8-A instructions all use 32 bits; contrary to RISC-V (and to the
``Thumb'' instructions on 32-bit Arm profiles), there are no compressed
16-bit instructions, thus little to gain there. This simplifies register
allocation, since general-purpose registers are interchangeable. The
local function call (\verb+bl+) uses a single 32-bit instruction; so does
the \verb+adr+ instruction, which loads the address of a function or a
static variable in a register. There is therefore no need to get the
address of \verb+decode_int()+ in a register, or to locate the block of
constants immediately after that function.

Most size optimization efforts on Armv8-A are spent trying to leverage
dual-operation instructions, such as memory accesses that also update
the index register, or arithmetic and logic instructions where an extra
shift or rotation can be applied to the second operand. In total,
Armv8-A offers fewer opportunities for size optimizations than RISC-V
and in particular x86\footnote{To some extent, this why I find assembly
programming on Armv8-A more pleasant than on x86.}.

\subsection{Missing Optimizations}

A few extra optimizations are possible on the implementations, but I
refrained from applying them because they would have a substantial cost
on performance, or, more importantly, would impact the implementation
clarity and make it more difficult to run unit tests and validate the
correctness of the implementation.

The virtual CPU runs a loop with a conditional jump over an exit flag,
and there is a separate macro that winds down the context and translates
the exit status back into the convention expected by the API function
caller. These steps could be merged into the \verb+_OK+/\verb+_FAIL+
code sequence, which would save a few bytes; but it would make it quite
harder to write unit tests that exercise individual bytecode sequences,
unless the unit test functions are made to use a different interpreter,
which would remove most of the point of running unit tests (such tests
have value precisely because they exercise parts of the implementation,
and not test-specific code).

Current \verb+_CALL+ implementation is generic and supports arbitrary
nesting levels of \verb+_CALL+ and \verb+_RET+. The actual
implementation only uses one level, and could use a dedicated register
to remember the caller's instruction pointer, instead of using the
native stack; this would make implementation of \verb+_CALL+ and
\verb+_RET+ a bit shorter, and would allow using the native stack
pointer to reference the interpreter context. However, that would
prevent some \emph{potential} improvements in the bytecode in which
two-level calls would be useful (e.g. to merge sub-sequences of code in
\verb+point_add_to_W+).

The current \verb+point_add_to_W+ was obtained by directly translating
the Renes-Costello-Batina algorithm into 43 three-instruction sequences,
then doing manual optimizations. It is unlikely that the current state
is optimal in that respect; further inspection of the operations would
probably reveal a few opportunities for saving some more bytes.

The current modulus ($p$ or $n$) is currently set with the \verb+_SMOD+
instruction, and read accesses as variable \verb+_M+ are supported by
some pointer-swapping in instruction decoding. This is done so because a
modulus is not only an integer but also includes the Montgomery
coefficient $2^{w-s}(\umod{-1/m}{2^s})$, which is located immediately
after the modulus in RAM. This special treatment could be avoided by
making variable \verb+_M+ a normal value, but expanding all values to an
internal format with an extra slot for the Montgomery coefficient; this
slot would simply follow along in value copies, and would be otherwise
ignored except when the value is used as the modulus in Montgomery
multiplications.

Presumably, unleashing an AI on our current code could make it somewhat
shorter, but readability and correctness validation would be lost in the
process. This validation is a crucial element for cryptographic
implementations; this is covered in the next section.

% ----------------------------------------------------------------------

\newpage
\section{Comprehensive Range Analysis}\label{sec:range}

The description of the ECDSA implementation in
section~\ref{sec:bytecode} highlighted that the correctness of the
implementation relies on the following properties:
\begin{itemize}

    \item All the additions and subtractions produce the correct
    integer results, i.e. on inputs that represent integers $x$ and $y$,
    addition returns an output that represents $x + y$, and subtraction
    similarly returns $x - y$. This is not prescriptive on the actual
    representation (in redundant representation, the same integer may
    accept multiple representations).

    \item Montgomery multiplications return, for inputs $x$ and $y$, and
    current modulus $m$, an output $z$ such that $\eqvm{z}{xy/R}{m}$.
    This is again not prescriptive on representations: it does not
    require a specific $z$ value since only modular equality is needed
    here.

    \item Integer canonicalization properly returns an integer in
    the $[0, m-1]$ range, for the modulus $m$. This property is required
    for proper bit enumeration over the scalars $u$ and $v$ for the double
    point multiplication $uG + vQ$, as well as for correctly handling the
    reinterpretation of $x_W$ into a scalar.

\end{itemize}
This section is about proving these properties for any possible inputs
to the ECDSA verification algorithm. This is not a \emph{complete}
formal proof of the entire implementation, in that it does not cover
some needed tests, such as verifying that the input $r$ and $s$ are in
the $[1, n-1]$ range, or the public key point is indeed a curve point.
Such properties are amenable to a manual, human-based verification by
simply looking at the code; contrary to rare mathematical edge cases,
notably related to carry propagation, their analysis should be simple.
You can add some explicit test vectors if you think that your developers
would manage to goof up something as simple as, for instance, verifying
that the input hash value length is between 32 and 64 bytes\footnote{My
unit test framework has such tests because I did goof it up at some
point.}.

The proving operation is done in a Python script called
\verb+rr_ecdsa_range.py+ and provided with the ECDSA implementations.

\subsection{Range Tracking}

The first two properties boil down to verifying that none of the
operations, which we expressed in section~\ref{sec:montymul} as
functions \verb+add()+, \verb+sub()+, \verb+longmul_s()+, etc, overflows
its expected container; we check this no-overflow property by doing a
range analysis which keeps track of possible values as ranges.
We define a \emph{range} as a pair of integers $(L, H)$ which mean that
the value is an element of the interval $[L, H]$ (note that both ends
are included). When $L = H$, the range is said to be \emph{exact} since
it pinpoints a single value.

In \verb+rr_ecdsa_range.py+, the \verb+ZRange+ class is defined; each
instance is immutable and represents a range. Elementary operations on
ranges are supported, such as additions, multiplications, etc; for
instance, when multiplying range $(L_1, H_1)$ with range $(L_2, H_2)$,
the output range is $(L_3, H_3)$ which is such that $L_3$ and $H_3$ are,
respectively, the minimum and the maximum in the quadruplet $\{ L_1 L_2,
L_1 H_2, H_1 L_2, H_1 H_2 \}$. We also need to support masking values to
their $t$ low bits: signed and unsigned $t$-bit masking of an integer
$x$ mean returning $\smod{x}{2^t}$ and $\umod{x}{2^t}$, respectively.
For ranges, the implementation of unsigned masking returns the full
range $[0, 2^t - 1]$ if the input $[L, H]$ is such that $H - L \geq 2^t$
(all low $t$-bit patterns are possible) or if $\umod{H}{2^t} <
\umod{L}{2^t}$ (which implies that both $0$ and $2^t - 1$ are possible,
thus encompassing the whole range, even if some of the values in the
range are not actually possible); otherwise, the relevant sub-range of
$[0, 2^t - 1]$ is returned. Signed masking uses unsigned masking with an
appropriate offset applied before and after the operation.

Building on \verb+ZRange+, the \verb+ZWord+ class represents a machine
word of $w$ bits. Each instance encapsulates two \verb+ZRange+
instances: one represents the possible range for the value stored in the
word, and the other is the container range, which is either $[-2^{w-1},
2^{w-1} - 1]$ (the word is interpreted as a signed integer) or $[0, 2^w
- 1]$ (the word is interpreted as an unsigned integer). Arithmetic
operations are implemented on the value range, and automatically raise
an exception if the output value range would not entirely fit within the
container range: this is how overflows on word operations are tracked.
Explicit ``wrapping'' operations are also provided, which truncate the
output as appropriate instead of raising an exception; this is used for
the computation of the Montgomery factors $f_i$, for which the
truncation, and reinterpretation of an unsigned value as a signed
integer, are part of the operation semantics.

The \verb+MLInt+ class builds on \verb+ZRange+ and \verb+ZWord+ to
implement a multi-limb redundant representation of an integer. Each
\verb+MLInt+ instance has three fields:
\begin{itemize}

    \item A reference to an overall context structure which contains
    parameters such as $k$ (number of limbs), $s$ (for radix $2^s$)
    or flags such as whether signed or unsigned normalization should
    follow a Montgomery multiplication.

    \item The integer limbs, as an array of $k$ instances of
    \verb+ZWord+ (with signed interpretation): we keep track of the
    individual ranges of all limbs separately.

    \item An overall integer range (as a \verb+ZRange+), which qualifies
    the represented big integer. This may be a strict sub-range of
    the possible values that could be inferred from the ranges of
    the individual limbs.

\end{itemize}
Again, addition, subtraction and negation are implemented, as well
as Montgomery multiplications. This last operation uses a modulus
as extra parameter, provided as a \verb+MLIntModulus+ instance, which
contains the modulus limbs --- as \verb+ZWord+ values, which are not
necessarily exact if we want to do range analysis with moduli which
are not known at compile-time --- and some extra parameters such as
the Montgomery coefficient $\umod{-1/m}{2^s}$.

All the \verb+ZRange+, \verb+ZWord+ and \verb+MLInt+ instances are
immutable; all mathematical operations return new instances with the
result. Using immutable instances increases memory allocation traffic
(hence garbage collection costs) but avoids all possible issues with
concurrent modifications to the same object.

Since \verb+ZWord+ instances automatically check for overflows, executing
the operations without any thrown exception implies a correct result
for all input values fitting in the provided input ranges.

\subsection{Application to ECDSA Verification}

We can use \verb+MLInt+ in a full ECDSA verification: we set the input
ranges for signature elements, public key coordinates, and hash value
to the full 256-bit range $[0, 2^{256} - 1]$, and then run the algorithm
to completion. This requires a representation of the ECDSA code as
a sequence of operations that the Python script can map to operations
on \verb+MLInt+ instances. But we already have such a representation:
the bytecode described in section~\ref{sec:bytecode}.

The \verb+rr_ecdsa_range.py+ thus includes a textual copy of the bytecode
(as found in the assembly implementations) coupled with a minimal parser
that turns it into an internal representation of the bytecode instructions;
a bytecode interpreter then executes the instruction just like the
virtual machine assembly implementations, except that it maintains the
state variables as \verb+MLInt+ instances. Two API functions are
provided:
\begin{itemize}

    \item \verb+test_ecdsa()+ uses exact ranges as inputs (signature,
    public key, hash value), i.e. ranges that pinpoint a single value,
    obtained from a test vector. All 492 Wycheproof test vectors are
    used. Since the input values are exact, all derived values in the
    bytecode execution are also exact, and this just runs the ECDSA
    verifier. Going through all the test vectors is meant to confirm
    that the bytecode interpreter in \verb+rr_ecdsa_range.py+ is
    properly implemented\footnote{Running all 492 tests can take a
    lot of time, more than one hour for some choices of $k$ and $s$,
    due to the inherent inefficiency of keeping track of ranges and
    allocating all these values.}.

    \item \verb+range_analysis()+ instead sets inputs to the full
    256-bit range. As the code unfolds, maximum possible ranges for all
    words and big integers are computed at all points of the execution;
    if the end of the computation is reached (the final \verb+_OK+ or
    \verb+_FAIL+) without triggering any overflow check, then this
    proves that no overflow is indeed possible in the implementation.
    Additionally, extra tests are hooked on calls to the
    \verb+canonicalize+ routine, to check that the obtained big integer
    value range is indeed $[0, m-1]$ for the relevant modulus $m$ ($n$
    or $p$, depending on the location of the call).

\end{itemize}

One non-trivial point about this range analysis must be explained: the
executed code sequence does not always run the same list of
instructions, since conditional execution is used through the three
provided ``skip'' instructions (\verb+_SKIPBITZ+, \verb+_SKIPNZ+ and
\verb+_SKIPNEG+). When the tested value is known with an exact range, or
at least with a range that is precise enough to yield the information
(e.g. for \verb+_SKIPNEG+, if the top limb range is entirely negative or
entirely non-negative), then the interpreter can apply the skip
semantics; otherwise, the skip operation is said to be \emph{unresolved}
and the interpreter must assume that both behaviours (skip and non-skip)
are possible, run \emph{both}, and merge the resulting states. Two
\verb+MLInt+ values are merged by pairwise merging the ranges of their
limbs and overall values, the merge of two ranges being the smallest
range that includes both source ranges. This process must handle a few
situations:
\begin{itemize}

    \item When the skip operation is over a \verb+_CALL+ instruction,
    the called function must be invoked with a fresh copy of the
    VM state, so that after the function execution, the output state
    resulting from the function invocation can be merged into the
    saved state.

    \item Unresolved skips over \verb+_OK+ or \verb+_FAIL+ are simply
    ignored, since then the no-skip branch exits the VM and contributes
    no state to merge.

    \item To handle conditional returns in functions (unresolved skips
    over \verb+_RET+), the call stack records not only the caller's
    instruction pointer, but also a slot for a ``merged state'' which
    corresponds to the output states of all conditional returns. This
    allows computing the overall output state of the function, for all
    possible ways the function may exit.

    \item For \verb+_SKIPNEG+, the two branches are executed over states
    that take into account the result of the test: the skipped instruction
    uses a state that includes the information that the top limb of the
    value in the accumulator was non-negative (the overall integer value
    range is also adjusted accordingly), while the saved state (for the
    ``instruction skipped'' branch) includes the information that the top
    limb of the value in the accumulator was negative. This is needed in
    order to prove that canonicalization works properly.

    \item If the skip instruction applies to \verb+_NEXT+,
    \verb+_FNEXT+, or another skip instruction, then this is explicitly
    detected and an exception is thrown. This case does not happen in
    the ECDSA verifier bytecode. Supporting these cases would complicate
    the range analysis, since the simple execute-and-merge process would
    no longer suffice.

\end{itemize}

\newenvironment{codedesc4}
    {\begin{list}
        {}
        {\setlength{\labelwidth}{3em}
         \setlength{\labelsep}{0.5em}
         \setlength{\itemsep}{1ex}
         \setlength{\leftmargin}{3.5em}
         \setlength{\rightmargin}{0em}
         \let\makelabel=\makecodedesclabel
        }
    }
    {\end{list}}

The \verb+test_ecdsa()+ and \verb+range_analysis()+ functions accept the
same parameters, the first three of which being mandatory:
\begin{codedesc4}

    \item[k] Number of limbs (integer).

    \item[s] Limb radix size, in bits (integer).

    \item[w] Machine word size, in bits (integer).

    \item[slf] If \verb+True+, then the Montgomery factors use signed
    interpretation (a \verb+longmul_s+ is used for multiplying $f_i$
    with modulus limb $m_j$, as shown in
    listing~\ref{alg:mmul-redundant}). For the integer canonicalization,
    as implemented in listing~\ref{alg:canonicalize}, to work properly,
    \verb+slf+ must be set to \verb+True+. Default value is \verb+False+.

    \item[slv] If \verb+True+, then signed normalization is applied
    at the end of Montgomery multiplication; default value is
    \verb+False+, which means that unsigned normalization is applied
    (as in listing~\ref{alg:mmul-redundant}).

    \item[slm] If \verb+True+, then the moduli ($p$ and $n$) are stored
    and used in signed-normalized representation; default value is
    \verb+False+, which implies that the moduli are unsigned-normalized.

    \item[spec] A set of flags, as an integer which is a bitwise \textsc{or}
    of individual flags (flag $j$ has numerical value $2^j$). When set,
    each flag adds a \verb+_NORM+ instruction at a specific place; these
    extra unsigned normalization instructions are needed in some cases,
    depending on the values of the other parameters (this will be
    detailed below). Default value is 0 (all flags cleared).

\end{codedesc4}
A sample usage in a Python command-line prompt is the following:
{\small\begin{verbatim}
  $ python3
  Python 3.12.3 (main, Mar  3 2026, 12:15:18) [GCC 13.3.0] on linux
  Type "help", "copyright", "credits" or "license" for more information.
  >>> from rr_ecdsa_range import *
  >>> range_analysis(5, 54, 64, slf=True)
  OK
  >>> test_ecdsa(5, 54, 64, slf=True)
  ................................ (etc)
\end{verbatim}}
The fact that \verb+range_analysis()+ prints ``OK'' and exits without
raising an exception shows that when $k = 5$, $s = 54$ (54-bit limbs),
$w = 64$ (64-bit limbs), and Montgomery factors $f_i$ are interpreted
as signed values, then the entire ECDSA implementation never overflows
for any input. Running \verb+test_ecdsa()+ furthermore shows that the
computation is correct enough for all 492 test vectors.

\subsection{Variants and Flags}\label{sec:range:variants}

The effect of the \verb+slf+, \verb+slv+, \verb+slm+ and \verb+spec+
parameters is best explained with examples. If we try the calls above
without \verb+slf=True+, then Montgomery factors are unsigned and the
calls fail. For instance:
{\small\begin{verbatim}
  >>> range_analysis(5, 54, 64)
  <current IP = 8 [line 14]: ST 5>
  (...)
      assert x.low >= 0 and x.high <= mod_N - 1
             ^^^^^^^^^^^^^^^^^^^^^^^^^^^^^^^^^^
  AssertionError
\end{verbatim}}
The designated instruction is the \verb+_ST+ that immediately follows
a call to \verb+canonicalize+ (with modulus $n$): the canonicalization
relies on Montgomery reduction to return a value in $[-\gamma n, +\gamma n]$
for some constant $\gamma < 1$, and this is ensured only if the Montgomery
factors are signed; otherwise, the overall reduction factor $f$ can become
greater than $R/2$, and on output of Montgomery reduction, one may obtain
an integer equal to or slightly greater than the modulus, which is not
correctly handled by the current implementation of \verb+canonicalize+,
which does only one conditional \emph{addition} of the modulus, but not
a conditional \emph{subtraction}. All our implementations thus use
\verb+slf=True+.

The $(k, s, w) = (5, 54, 64)$ setting with \verb+slf=True+ is what we
implement in the 64-bit C code (\verb+gen64+) and all three 64-bit
assembly implementations (\verb+amd64+, \verb+arm64+ and \verb+rv64+).

If we set \verb+slm=True+, then range analysis still passes successfully,
but unit tests fail:
{\small\begin{verbatim}
  >>> range_analysis(5, 54, 64, slf=True, slm=True)
  OK
  >>> test_ecdsa(5, 54, 64, slf=True, slm=True)
  (...)
      assert status == 'F'
             ^^^^^^^^^^^^^
  AssertionError
\end{verbatim}}
The effect of setting \verb+slm+ to \verb+True+ is that the moduli are
internally signed-normalized instead of unsigned-normalized. For
instance, in signed-normalized format, the two low limbs of $p$ are $-1$
and $2^{42}$, instead of $2^{54}-1$ and $2^{42}-1$, respectively. Signed
normalization of the modulus has the effect that when adding $f_i m$ to
the intermediate result in Montgomery representation, the integer
products $f_i m_j$ are smaller in absolute value, which can help with
avoiding an integer overflow in the limbs. The effect is usually
marginal, but with some combinations of $k$, $s$ and $w$, it can
potentially help succeed through a sequence of operations that would
otherwise potentially trigger an overflow. Unfortunately, using a
signed-normalized modulus breaks the \verb+invert_mod+ implementation
(listing~\ref{alg:invert_mod}) because the FLT code needs the exponent
$m-2$ in canonicalized (unsigned-normalized) format and expects to
obtain it by simply subtracting 2 from the low limb, which does not work
if $m$ is signed-normalized. The solution is to add an extra \verb+_NORM+
instruction to unsigned-normalize the exponent, which we can do by setting
the special flag 0 through the \verb+spec+ parameter:
{\small\begin{verbatim}
  >>> test_ecdsa(5, 54, 64, slf=True, slm=True, spec=1)
  ................................ (etc)
\end{verbatim}}
This enlarges the bytecode by 1 byte.

If we try to make a 32-bit version, we can find that with only
\verb+slf=True+, range analysis finds no workable solution with $k\leq 12$.
For instance:
{\small\begin{verbatim}
  >>> range_analysis(12, 22, 32, slf=True)
  <current IP = 198 [line 217]: MUL 23>
  <current IP = 54 [line 50]: SKIPBITZ 4>
  (...)
  Exception: value range [-2273312274, 1056964453] exceeds
  container range [-2147483648, 2147483647]
\end{verbatim}}
This indicates that a limb potentially overflows. The printed reference
to the offending instruction points to a \verb+_MUL+ opcode which is
near the end of the \verb+point_add_to_W+ routine. We can then try to
put a \verb+_NORM+ call right before this multiplication: unsigned
normalization makes individual limbs smaller, and the extra \verb+_NORM+
will do that on the accumulator, which is one of the operands of the
\verb+_MUL+. In the bytecode embedded in \verb+rr_ecdsa_range.py+,
we placed such a \verb+_NORM+ instruction, enabled with special flag 2
(i.e. by setting \verb+spec+ to 4):
{\small\begin{verbatim}
  >>> range_analysis(12, 22, 32, slf=True, spec=4)
  OK
  >>> test_ecdsa(12, 22, 32, slf=True, spec=4)
  ................................ (etc)
\end{verbatim}}
This extra \verb+_NORM+ enlarges the bytecode by 1 byte. This is the
setting used in the alternate x86 assembly implementation
(\verb+amd64alt+).

Another option for a 12-limb 32-bit implementation is to use signed
normalization at the end of each Montgomery multiplication (with
\verb+slv=True+). Signed normalization makes output limbs smaller (in
absolute value) and thus should allow more successive additions and
subtractions before the next multiplication. However, this would
predictably break canonicalization: the \verb+canonicalize+ routine
(listing~\ref{alg:canonicalize}) uses Montgomery reduction
(\verb+_MRED+) then tests the sign of the output with \verb+_SKIPNEG+,
under the assumption that the top limb will contain the value sign,
which is true for unsigned-normalized values, but not for
signed-normalized values. When \verb+slv=True+, we thus need an extra
\verb+_NORM+ between the \verb+_MRED+ and the \verb+_SKIPNEG+, which in
\verb+rr_ecdsa_range.py+ is obtained by setting the special flag 1
(value 2 into \verb+spec+):
{\small\begin{verbatim}
  >>> range_analysis(12, 22, 32, slf=True, slv=True, spec=2)
  OK
  >>> test_ecdsa(12, 22, 32, slf=True, slv=True, spec=2)
  ................................ (etc)
\end{verbatim}}
This specific setup is what is used in the 32-bit portable C
implementation (\verb+gen32+). In general, it is slightly less desirable
than the previous option (with \verb+slv=False+ and an extra
\verb+_NORM+ in \verb+point_add_to_W+) because it means that the
compiled code will contain instructions for \emph{both} signed and
unsigned normalization, instead of just unsigned normalization.

An intriguing third option for a 12-limb 32-bit implementation is to
use a signed modulus (\verb+slm=True+, and special flag 0 set). It fails
range analysis, but test vectors pass successfully, which can only mean
one of two things: either the range analysis is not tight enough, and
its inherent approximations lead it to reject a valid solution; or the
value overflow is in fact possible, but rare, and the test vectors do
not exercise the faulty condition. Avoiding rare edge case is exactly
why we are doing an extensive range analysis, and we must therefore
reject that option until an actual correctness proof can be established.

% ----------------------------------------------------------------------

\newpage
\section{Conclusion}\label{sec:conclusion}

In this paper, we described an implementation technique for Montgomery
multiplications that avoids most carry propagations and should be
desirable for computations on platforms for which carry propagation is
expensive, in particular RISC-V CPUs. The use of redundant
representations with signed limbs also opens the possibility of
performing unreduced limb-wise additions and subtractions, which may
yield non-negligible performance benefits in a speed-oriented
implementation of relatively complex code sequences, as is common in
schemes using elliptic curves. This should be applicable, in particular,
to post-quantum schemes based on isogeny graphs.

An additional potential benefit of redundant representations is, for
moduli with a special format, selection of an alternate representation
that would allow more efficient operations; for instance, the parameters
proposed for the 128-bit security level of FESTA\cite{BasMaiPop2023}, a
post-quantum isogeny-based public key encryption scheme, use a modulus
$m$ such that $m+1$ is a multiple of $2^{632}$; in a classic unsigned
and non-redundant representation, this leads to a sequence of low limbs
with a large value. However, in a signed redundant representation, we
could set the lowest limb to $-1$, and then quite a few limbs beyond it
will have value $0$. Multiplication by $-1$ is efficient, multiplication
by $0$ even more so. The extent of such optimizations depends on the
modulus value, but might lead to non-negligible speed improvements. Best
case seems to be moduli whose binary representation includes a very long
sequence of ones, so that signed limbs may turn all these into zeros.

We illustrated the technique with an example of questionable usefulness
(a codegolfed implementation of ECDSA signature verification over NIST
curve P-256), but appropriate for highlighting how a comprehensive range
analysis can be performed, as is required in order to use the unreduced
limb-wise additions and subtractions. Python-based tooling can in
particular pinpoint where overflows may occur, and help decide where
extra normalization steps need to be placed. The proof of impossibility
of overflows, combined with the removal of most carry propagation,
should help in making correct implementations.

In the implementation process, we also estimated how reliable and useful
the current crop of AIs can be for implementing cryptographic
primitives. Verdict is: they produce results quickly, but be very
cautious, they tend to focus on test vectors, which cannot guarantee
correctness. You really should ask the AI to produce a companion
verifiable proof of correctness.

% ----------------------------------------------------------------------

\section*{Acknowledgements}\label{sec:acknowledge}

Many thanks to Steve Weis, who provided access to the AI.

% ----------------------------------------------------------------------

\newpage
\begin{thebibliography}{20}

% Ideally we would do that only for URLs, but I don't know how to do
% that in a non-painful way.
\RaggedRight

%\bibitem{AbdHwaKanSpr2022}
%A.~Abdulrahman, V.~Hwang, M.~Kannwischer and A.~Sprenkels,
%\emph{Faster Kyber and Dilithium on the Cortex-M4},
%Applied Cryptography and Network Security -- ACNS~2022,
%Lecture Notes In Computer Science, vol.~13269, pp.~853-871, 2022.

%\bibitem{AAPCS}
%\emph{Procedure Call Standard for the Arm® Architecture},
%Arm Limited, version 2024Q3, 2024.

%\bibitem{ChoYooSeo2024}
%J.~Choi, S.~Yoon and S.C.~Seo,
%\emph{Optimized Falcon Verify on Cortex-M4 for Post-Quantum secure UAV 
%communications},
%ICT Express, Elsevier, 2024,\\
%\url{https://doi.org/10.1016/j.icte.2024.11.002}

%\bibitem{DuzFieFis2024}
%S.~Düzlü, R.~Fiedler and M.~Fischlin,
%\emph{BUFFing FALCON without Increasing the Signature Size},\\
%\url{https://eprint.iacr.org/2024/710}

%\bibitem{GajJanKin2024}
%P.~Gajland, J.~Janneck and E.~Kiltz,
%\emph{A Closer Look at Falcon},\\
%\url{https://eprint.iacr.org/2024/1769}

%\bibitem{Hwa2024}
%V.~Hwang,
%\emph{Formal Verification of Emulated Floating-Point Arithmetic in Falcon},\\
%Advances in Information and Computer Security: 19th International Workshop
%on Security -- IWSEC 2024, Lecture Notes in Computer Science,
%vol.~14977, pp.~125-141, 2024.

%\bibitem{ieee754}
%\emph{IEEE Standard for Binary Floating-Point Arithmetic}, 1985,\\
%\url{https://doi.org/10.1109%2FIEEESTD.1985.82928}

%\bibitem{Gol1991}
%D.~Goldberg,
%\emph{What Every Computer Scientist Should Know About Floating-Point
%Arithmetic},
%ACM Computing Surveys (CSUR), vol.~23, issue~1, pp.~5-48, 1991.

%\bibitem{KanRijSchSto2019}
%M.~Kannwischer, J.~Rijneveld, P.~Schwabe and K.~Stoffelen,
%\emph{pqm4: Testing and Benchmarking NIST PQC on ARM Cortex-M4},\\
%\url{https://eprint.iacr.org/2019/844}

%\bibitem{NISTPQC}
%\emph{Post-Quantum Cryptography},
%National Institute of Standard and Technology,\\
%\url{https://www.nist.gov/pqcrypto}

%\bibitem{Por2019}
%T.~Pornin,
%\emph{New Efficient, Constant-Time Implementations of Falcon},\\
%\url{https://eprint.iacr.org/2019/893}

%\bibitem{Por2020-1}
%T.~Pornin,
%\emph{Efficient Elliptic Curve Operations On Microcontrollers With
%Finite Field Extensions},\\
%\url{https://eprint.iacr.org/2020/009}

%\bibitem{Por2023-1}
%T.~Pornin,
%\emph{Improved Key Pair Generation for Falcon, BAT and Hawk},\\
%\url{https://eprint.iacr.org/2023/290}

%\bibitem{XKCP}
%G.~Bertoni, J.~Daemen, S.~Hoffert, M.~Peeters, G.~Van~Assche and
%R.~Van~Keer,
%\emph{eXtended Keccak Code Package},\\
%\url{https://github.com/XKCP/XKCP}

%\bibitem{AarAra2022}
%M.~Aardal and D.~Aranha,
%\emph{2DT-GLS: Faster and exception-free scalar multiplication in the
%GLS254 binary curve},\\
%\url{https://eprint.iacr.org/2022/748}

%\bibitem{AmbBosFayJoyLocMur2017}
%C.~Ambrose, J.~Bos, B.~Fay, M.~Joye, M.~Lochter and B.~Murray,
%\emph{Differential Attacks on Deterministic Signatures},\\
%\url{https://eprint.iacr.org/2017/975}

\bibitem{X962}
Accredited Standards Committee X9,
\emph{The Elliptic Curve Digital Signature Algorithm (ECDSA)},
American National Standards Institute, ANS X9.62, 2005.

\bibitem{X963}
Accredited Standards Committee X9,
\emph{Key Agreement and Key Transport Using Elliptic Curve Cryptography},
American National Standards Institute, ANS X9.63, 2017.

%\bibitem{AntBroGalLamStrVan2005}
%A.~Antipa, D.~Brown, R.~Gallant, R.~Lambert, R.~Struik and S.~Vanstone,
%\emph{Accelerated Verification of ECDSA signatures},
%Selected Areas in Cryptography - SAC 2005, Lecture Notes in Computer
%Science, vol.~3897, pp.~307-318, 2005.

%\bibitem{AraLopHan2010}
%D.~Aranha, J.~López and D.~Hankerson,
%\emph{Efficient Software Implementation of Binary Field Arithmetic Using
%Vector Instruction Sets},
%Progress in Cryptology - LATINCRYPT 2010,
%Lecture Notes in Computer Science, vol.~6212, pp.~144-161, 2010.

%\bibitem{RistrettoWeb}
%T.~Arcieri, I.~Lovecruft and H.~de Valence,
%\emph{The Ristretto Group},\\
%\url{https://ristretto.group/}

%\bibitem{Atk1992}
%A.~Atkin,
%\emph{Probabilistic primality testing (summary by F. Morain)}, Technical Report 1779, INRIA, 1992,\\
%\url{http://algo.inria.fr/seminars/sem91-92/atkin.pdf}

\bibitem{BasMaiPop2023}
A.~Basso, L.~Maino and G.~Pope,
\emph{FESTA: Fast Encryption from Supersingular Torsion Attacks},\\
\url{https://ia.cr/2023/660}

\bibitem{Ber2006}
D.~Bernstein,
\emph{Curve25519: new Diffie-Hellman speed records},
Public Key Cryptography - PKC 2006,
Lecture Notes in Computer Science, vol.~3958, pp.~207-228, 2006.

\bibitem{BerDuiLanSchYan2012}
D.~Bernstein, N.~Duif, T.~Lange, P.~Schwabe and B.-Y.~Yang,
\emph{High-speed high-security signatures},
Journal of Cryptographic Engineering, vol.~2, issue~2, pp.~77-89, 2012.

%\bibitem{eBACS}
%D.~Bernstein and T.~Lange,
%\emph{eBACS: ECRYPT Benchmarking of Cryptographic Systems},\\
%\url{https://bench.cr.yp.to} (accessed 4 August 2022).

%\bibitem{BerLanRez2008}
%D.~Bernstein, T.~Lange and R.~Rezaeian Farashahi,
%\emph{Binary Edwards Curves},
%Cryptographic Hardware and Embedded Systems - CHES 2008,
%Lecture Notes in Computer Science, vol.~5154, pp.~244-265, 2008.

%\bibitem{BerHamKraLan2013}
%D.~Bernstein, M.~Hamburg, A.~Krasnova and T.~Lange,
%\emph{Elligator: elliptic-curve points indistinguishable from uniform
%random strings},
%Proceedings of the 2013 ACM SIGSAC Conference on Computer and Communications
%Security, 2013,\\
%\url{https://doi.org/10.1145/2508859.2516734}

%\bibitem{BerYan2019}
%D.~Bernstein and B.-Y.~Yang,
%\emph{Fast constant-time gcd computation and modular inversion},\\
%\url{https://gcd.cr.yp.to/papers.html#safegcd}

%\bibitem{BilJoy2003}
%O.~Billet and M.~Joye,
%\emph{The Jacobi model of an elliptic curve and side-channel analysis},
%AAECC-15, Lecture Notes in Computer Science, vol.~2643, pp.~34-42, 2003.

%\bibitem{Bod2007}
%M.~Bodrato,
%\emph{Towards Optimal Toom-Cook Multiplication for Univariate and
%Multivariate Polynomials in Characteristic 2 and 0},
%International Workshop on the Arithmetic of Finite Fields - WAIFI 2007,
%Lecture Notes in Computer Science, vol.~4547, pp.~116-133, 2007.

%\bibitem{BreGauThoZim2008}
%R.~Brent, P.~Gaudry, E.~Thomé and P.~Zimmerman,
%\emph{Faster multiplication in GF(2)[x]},
%Algorithmic Number Theory - ANTS-VIII, Lecture Notes in Computer Science,
%vol.~5011, pp.~153-166, 2008.

%\bibitem{BriCorIcaMadRanTib2010}
%É.~Brier, J.-S.~Coron, T.~Icart, D.~Madore, H.~Randriam and M.~Tibouchi,
%\emph{Efficient Indifferentiable Hashing into Ordinary Elliptic Curves},
%Advances in Cryptology - CRYPTO 2010, Lecture Notes in Computer Science,
%vol.~6223, pp.~237-254, 2010.

\bibitem{Bri1983}
E.~Brickell,
\emph{A Fast Modular Multiplication Algorithm with Application to Two Key
Cryptography},
Advances in Cryptology - CRYPTO 1982, pp.~51-60, 1983

%\bibitem{BruCurHof1993}
%H.~Brunner, A.~Curriger and M.~Hofstetter,
%\emph{On computing multiplicative inverses in GF($2^m$)},
%IEEE Transactions on Computers, vol.~42, issue~8, pp.~1010-1015, 1993.

%\bibitem{CanKal1991}
%D.~Cantor and E.~Kaltofen,
%\emph{On fast multiplication of polynomials over arbitrary algebras},
%Acta Informatica, vol.~28, issue~7, pp.~693-701, 1991.

%\bibitem{SEC2v2}
%Certicom Research,
%\emph{SEC 2: Recommended Elliptic Curve Domain Parameters},
%version~2.0, 2010,\\
%\url{https://www.secg.org/sec2-v2.pdf}

%\bibitem{Cha2022}
%K.~Chalkias,
%\emph{ed25519-unsafe-libs},\\
%\url{https://github.com/MystenLabs/ed25519-unsafe-libs}

%\bibitem{ChaGarNik2020}
%K.~Chalkias, F.~Garillot and V.~Nikolaenko,
%\emph{Taming the Many EdDSAs},
%Security Standardisation Research - SSR 2020, Lecture Notes in
%Computer Science, vol.~12529, pp.~67-90, 2020.

%\bibitem{ChuChu1986}
%D.~Chudnovsky and G.~Chudnovsky,
%\emph{Sequences of numbers generated by addition
%in formal groups and new primality and factorization tests},
%Advances in Applied Mathematics, vol.~7, issue~4, pp.~385–434, 1986.

%\bibitem{CohFre2006}
%H.~Cohen and G.~Frey (editors),
%\emph{Handbook of Elliptic and Hyperelliptic Curve Cryptography},
%Chapman and Hall/CRC, 2006.

\bibitem{Col1726}
J.~Colson,
\emph{A Short Account of Negativo-Affirmative Arithmetick},
Philosophical Transactions (1683-1775), vol.~34, pp.~161-173, 1726.

%\bibitem{FrostDraft}
%D.~Connolly, C.~Komlo, I.~Goldberg and C.~Wood,
%\emph{Two-Round Threshold Schnorr Signatures with FROST},\\
%\url{https://datatracker.ietf.org/doc/html/draft-irtf-cfrg-frost-08}

%\bibitem{CreJac2019}
%C.~Cremers and D.~Jackson,
%\emph{Prime, Order Please! Revisiting Small Subgroup and Invalid Curve
%Attacks on Protocols using Diffie-Hellman},
%IEEE 32nd Computer Security Foundations Symposium (CSF), 2019.

%\bibitem{dalek-crypto}
%I.~Lovecruft and H.~de Valence,
%\emph{Dalek cryptography},\\
%\url{https://github.com/dalek-cryptography}

\bibitem{DifHel1976}
W.~Diffie and M.~Hellman,
\emph{New directions in cryptography},
IEEE Transactions on Information Theory, vol.~22, issue~6, pp.~644-654, 1976.

%\bibitem{Duq2007}
%S.~Duquesne,
%\emph{Improving the arithmetic of elliptic curves in the Jacobi model},
%Information Processing Letters, vol.~104, issue~3, pp.~101–105, 2007.

%\bibitem{FisPat1974}
%M.~Fischer and M.~Paterson,
%\emph{String-Matching and Other Products},
%MAC Technical Memorandum 41, Massachusetts Institute of Technology, 1974.

%\bibitem{Fog2023}
%A.~Fog,
%\emph{The microarchitecture of Intel, AMD, and VIA CPUs},\\
%\url{https://www.agner.org/optimize/microarchitecture.pdf}

%\bibitem{GalLamVan2001}
%R.~Gallant, J.~Lambert and S.~Vanstone,
%\emph{Faster Point Multiplication on Elliptic Curves with Efficient
%Endomorphisms},
%Advances in Cryptology - CRYPTO 2001, Lecture Notes in Computer Science,
%vol.~2139, pp.~190-200, 2001.

%\bibitem{GalLinSco2009}
%S.~Galbraith, X.~Lin and M.~Scott,
%\emph{Endomorphisms for Faster Elliptic Curve Cryptography on a Large
%Class of Curves},
%Advances in Cryptology - EUROCRYPT 2009, Lecture Notes in Compute Science,
%vol.~5479, pp.~518-535, 2009.

%\bibitem{GauHesSma2002}
%P.~Gaudry, F.~Hess and N.~Smart,
%\emph{Constructive and destructive facets of Weil descent on elliptic curves},
%Journal of Cryptology, vol.~15, issue~1, pp.~19-46, 2002.

%\bibitem{Ham2009}
%M.~Hamburg,
%\emph{Accelerating AES with Vector Permute Instructions},
%Cryptographic Hardware and Embedded Systems - CHES 2009,
%Lecture Notes in Computer Science, vol.~5747, pp.~18-32, 2009.

%\bibitem{Ham2015-1}
%M.~Hamburg,
%\emph{Decaf: Eliminating cofactors through point compression},
%Advances in Cryptology - CRYPTO 2015, Lecture Notes in Computer Science,
%vol.~9215, pp.~705-723, 2015.

%\bibitem{Ham2015-2}
%M.~Hamburg,
%\emph{Ed448-Goldilocks, a new elliptic curve},\\
%\url{https://eprint.iacr.org/2015/625}

%\bibitem{HanKarMen2009}
%D.~Hankerson, K.~Karabina and A.~Menezes,
%\emph{Analyzing the Galbraith-Lin-Scott Point Multiplication Method for
%Elliptic Curves over Binary Fields},
%IEEE Transactions on Computers, vol.~58, issue~10, pp.~1411-1420, 2009.

%\bibitem{HanMenVan2004}
%D.~Hankerson, A.~Menezes and S.~Vanstone,
%\emph{Guide to Elliptic Curve Cryptography},
%Springer/Sci-Tech/Trade, 2004.

%\bibitem{HisWonCarDaw2008}
%H.~Hisil, K.~Wong, G.~Carter and E.~Dawson,
%\emph{Twisted Edwards Curves Revisited},
%Advances in Cryptology - ASIACRYPT 2008, Lecture Notes in Computer Science,
%vol.~5350, pp.~326-343, 2008.

%\bibitem{HisWonCarDaw2009}
%H.~Hisil, K.~Wong, G.~Carter and E.~Dawson,
%\emph{Jacobi Quartic Curves Revisited},
%Information Security and Privacy - ACISP 2009, Lecture Notes in Computer
%Science, vol.~5594, pp.~452-468, 2009.

\bibitem{Hon1989}
J.~Honkala,
\emph{On Number Systems with Negative Digits},
Annales Academiæ Scientiarum Fennicæ, Series A. I. Mathematica,
vol.~14, issue~1, pp.149-156, 1989.

%\bibitem{ItoTsu1988}
%T.~Itoh and S.~Tsujii,
%\emph{A Fast Algorithm for Computing Multiplicative Inverses in
%$\GF(2^m)$ Using Normal Bases},
%Information and Computation, vol.~78, pp.~171-177, 1988.

%\bibitem{Jac1829}
%C.~G.~J.~Jacobi,
%\emph{Fundamenta nova theoriae functionum ellipticarum},
%Sumtibus fratrum, 1829.

%\bibitem{EdDSArfc8032}
%S.~Josefsson and I.~Liusvaara,
%\emph{Edwards-Curve Digital Signature Algorithm (EdDSA)},\\
%\url{https://tools.ietf.org/html/rfc8032}

%\bibitem{KarOfm1962}
%А.~Карацуба and Ю.~Офман,
%\emph{Умножение Многозначных Чисел На Автоматах},
%Доклады Академи и наук СССР, vol.~145, issue~2, pp.~293-294, 1962.
%{\fontencoding{T2A}\fontfamily{NimbusSerif}\selectfont А.~Карацуба}
%and {\fontencoding{T2A}\fontfamily{NimbusSerif}\selectfont Ю.~Офман,
%\emph{Умножение Многозначных Чисел На Автоматах},
%Доклады Академи и наук СССР}, vol.~145, pp.~293-294, 1962.
%A.~Karatsuba and Y.~Ofman,
%\emph{Multiplication of Many-Digital Numbers by Automatic Computers},
%Proceedings of the USSR Academy of Sciences, vol.~145, pp.~293-294, 1962.

%\bibitem{KimKim2007}
%K.~Kim and S.~Kim,
%\emph{A New Method for Speeding Up Arithmetic on Elliptic Curves over
%Binary Fields},\\
%\url{https://eprint.iacr.org/2007/181}

%\bibitem{Knu1969}
%D.~Knuth,
%\emph{The Art of Computer Programming, volume~2: Seminumerical Algorithms},
%Addison-Wesley, 1969.

%\bibitem{KomGol2020}
%C.~Komlo and I.~Goldberg,
%\emph{FROST: Flexible Round-Optimized Schnorr Threshold Signatures},\\
%\url{https://eprint.iacr.org/2020/852}

%\bibitem{rfc7748}
%A.~Langley, M.~Hamburg and S.~Turner,
%\emph{Elliptic Curves for Security},\\
%\url{https://www.rfc-editor.org/rfc/rfc7748}

\bibitem{Lean}
{\fontspec{texgyreadventor-regular.otf}\selectfont{L\kern -0.1em\reflectbox{E}\kern -0.13em\raisebox{\depth}{\scalebox{1}[-1]{A}}\kern -0.13em{N}}} \emph{Programming Language},\\
\url{https://lean-lang.org/}

%\bibitem{Len2019}
%E.~Lenngren,
%\emph{X25519 for ARM AArch64},\\
%\url{https://github.com/Emill/X25519-AArch64}

%\bibitem{Len2020}
%E.~Lenngren,
%\emph{X25519 for ARM Cortex-M4 and other ARM processors},\\
%\url{https://github.com/Emill/X25519-Cortex-M4}

%\bibitem{LopDah1999}
%J.~López and R.~Dahab,
%\emph{Fast Multiplication on Elliptic Curve Over $\GF(2^m)$ without
%precomputation},
%Cryptographic Hardware and Embedded Systems - CHES'99,
%Lecture Notes in Computer Science, vol.~1717, pp.~316-327, 1999.

\bibitem{MenOorVan1996}
A.~Menezes, P.~van Oorschot and S.~Vanstone,
\emph{Handbook of Applied Cryptography},
CRC Press, 1996,\\
\url{https://cacr.uwaterloo.ca/hac/}

\bibitem{Mon1985}
P.~Montgomery,
\emph{Modular Multiplication Without Trial Division},
Mathematics of Computation, vol.~44, issue~170, pp.~519-521, 1985.

%\bibitem{MooShu2011}
%D.~Moody and D.~Shumow,
%\emph{Analogues of Vélu's Formulas for Isogenies on Alternate Models of
%Elliptic Curves},
%\url{https://eprint.iacr.org/2011/430}

%\bibitem{NacSte2001}
%D.~Naccache and J.~Stern,
%\emph{Signing on a Postcard},
%Financial Cryptography - FC 2000, Lecture Notes in Computer Science,
%vol.~1962, pp.~121-135, 2001.

%\bibitem{NatSar2020}
%K.~Nath and P.~Sarkar,
%\emph{Efficient 4-way Vectorizations of the Montgomery Ladder},\\
%\url{https://eprint.iacr.org/2020/378}

%\bibitem{NevSmaWar2009}
%G.~Neven, N.~P.~Smart and B.Warinschi,
%\emph{Hash function requirements for Schnorr signatures},
%Journal of Mathematical Cryptology, vol.~3, issue 1, pp.~69-87, 2009.

%\bibitem{Fips180}
%Information Technology Laboratory,
%\emph{Secure Hash Standard (SHS)},
%National Institute of Standard and Technology, FIPS~180-4, 2015.

\bibitem{Fips186}
Information Technology Laboratory,
\emph{Digital Signature Standard (DSS)},
National Institute of Standard and Technology, FIPS~186-5, 2023.

%\bibitem{Fips202}
%Information Technology Laboratory,
%\emph{SHA-3 Standard: Permutation-Based Hash and Extendable-Output
%Functions},
%National Institute of Standard and Technology, FIPS~202, 2015.

\bibitem{Fips203}
Information Technology Laboratory,
\emph{Module-Lattice-Based Key-Encapsulation Mechanism Standard},
National Institute of Standard and Technology, FIPS~203, 2024.

\bibitem{Fips204}
Information Technology Laboratory,
\emph{Module-Lattice-Based Digital Signature Standard},
National Institute of Standard and Technology, FIPS~204, 2024.

%\bibitem{BLAKE3}
%J.~O'Connor, J.-P.~Aumasson, S.~Neves and Z.~Wilcox-O'Hearn,
%\emph{BLAKE3},\\
%\url{https://github.com/BLAKE3-team/BLAKE3}

%\bibitem{OliLopAraRod2014}
%T.~Oliveira, J.~López-Hernández, D.~Aranha and F.~Rodríguez-Henríquez,
%\emph{Two is the fastest prime: lambda coordinates for binary elliptic
%curves},
%Journal of Cryptographic Engineering, vol.~4, issue~1, pp.~3-17, 2014.

%\bibitem{OliLopAraRod2016}
%T.~Oliveira, J.~López-Hernández, D.~Aranha and F.~Rodríguez-Henríquez,
%\emph{Improving the performance of the GLS254},
%Presented at the CHES 2016 Rump Session.

%\bibitem{PetQui2012}
%C.~Petit and J.-J.~Quisquater,
%\emph{On Polynomial Systems Arising from a Weil Descent},
%Advances in Cryptology - ASIACRYPT 2012,
%Lecture Notes in Computer Science, vol.~7658, pp.~451-466, 2012.

\bibitem{PhaKor1994}
D.~Phatak and I.~Koren,
\emph{Hybrid Signed-Digit Number Systems: A Unified Framework for
Redundant Number Representations With Bounded Carry Propagation Chains},
IEEE Transactions on Computers, vol.~43, issue~8, pp.~880-891, 1994.

%\bibitem{PodSomSchLocRos2018}
%D.~Poddebniak, J.~Somorovsky, S.~Schinzel, M.~Lochter and P.~Rösler,
%\emph{Attacking Deterministic Signature Schemes Using Fault Attacks},
%2018 IEEE European Symposium on Security and Privacy (EuroS\&P),
%pp.~338-352, 2018.

%\bibitem{DeterministicECDSArfc}
%T.~Pornin,
%\emph{Deterministic Usage of the Digital Signature Algorithm (DSA) and
%Elliptic Curve Digital Signature Algorithm (ECDSA)},\\
%\url{https://tools.ietf.org/html/rfc6979}

%\bibitem{BearSSL}
%T.~Pornin,
%\emph{BearSSL: a smaller SSL/TLS library},\\
%\url{https://www.bearssl.org/}

%\bibitem{Por2020-2}
%T.~Pornin,
%\emph{Optimized Lattice Basis Reduction In Dimension 2, and Fast Schnorr
%and EdDSA Signature Verification},\\
%\url{https://eprint.iacr.org/2020/454}

%\bibitem{Por2020-3}
%T.~Pornin,
%\emph{Optimized Binary GCD for Modular Inversion},\\
%\url{https://eprint.iacr.org/2020/972}

%\bibitem{Por2020-5}
%T.~Pornin,
%\emph{Double-Odd Elliptic Curves},\\
%\url{https://eprint.iacr.org/2020/1558}

%\bibitem{Por2022-4}
%T.~Pornin,
%\emph{Double-Odd Jacobi Quartic},\\
%\url{https://eprint.iacr.org/2022/1052}

%\bibitem{Por2022-6}
%T.~Pornin,
%\emph{Efficient and Complete Formulas for Binary Curves},\\
%\url{https://eprint.iacr.org/2022/1325}

%\bibitem{Por2023-3}
%T.~Pornin,
%\emph{On Multiplications with Unsaturated Limbs},\\
%\url{https://research.nccgroup.com/2023/09/18/on-multiplications-with-unsaturated-limbs/}

\bibitem{FalconRound3}
T.~Prest, P.-A.~Fouque, J.~Hoffstein, P.~Kirchner,
V.~Lyubashevsky, T.~Pornin, T.~Ricosset, G.~Seiler, W.~Whyte and Z.~Zhang,
\emph{Falcon},
Technical report, National Institute of Standards and Technology, 2019,
\url{https://csrc.nist.gov/Projects/post-quantum-cryptography/post-quantum-cryptography-standardization/round-3-submissions}

\bibitem{RenCosBat2016}
J.~Renes, C.~Costello and L.~Batina,
\emph{Complete Addition Formulas for Prime Order Elliptic Curves},
Advances in Cryptology - EUROCRYPT 2016, Lecture Notes in Compute Science,
vol.~9665, pp.~403-428, 2016.\\
\url{https://ia.cr/2015/1060}

\bibitem{RivShaAdl1978}
R.~Rivest, A.~Shamir and L.~Adleman,
\emph{A Method for Obtaining Digital Signatures and Public-Key Cryptosystems},
Communications of the ACM, vol.~21, issue~2, pp.~120-126, 1978.

%\bibitem{BLAKE2}
%M.-J.~Saarinen and J.-P.~Aumasson,
%\emph{The BLAKE2 Cryptographic Hash and Message Authentication Code (MAC)}\\
%\url{https://datatracker.ietf.org/doc/html/rfc7693}

%\bibitem{Sch1989}
%C.~Schnorr,
%\emph{Efficient identification and signatures for smart cards},
%Advances in Cryptology - CRYPTO '89, Lecture Notes in Computer Science,
%vol.~435, pp.~239-252, 1990.

%\bibitem{ShaWoe2006}
%A.~Shallue and C.~van de Woestijne,
%\emph{Construction of rational points on elliptic curves over finite
%fields},
%Algorithm Number Theory Symposium - ANTS 2006, Lecture Notes in Computer
%Science, vol.~4076, pp.~510-524, 2006.

%\bibitem{Sil2009}
%J.~Silverman,
%\emph{The Arithmetic of Elliptic Curves} (2nd edition),
%Springer New York, NY, 2009.

%\bibitem{Sol2000}
%J.~Solinas,
%\emph{Efficient Arithmetic on Koblitz Curves},
%Designs, Codes and Cryptography, vol.~19, issue~2-3, pp.~195-249, 2000.

%\bibitem{Str1964}
%E.~Straus,
%\emph{Addition chains of vectors (problem 5125)},
%American Mathematical Monthly, vol.~70, pp.~806-808, 1964.

%\bibitem{Too1963}
%А.~Тоом,
%\emph{О сложности схемы из функциональных злементов, реализирующей
%умножение целых чисел},
%Доклады Академи и наук СССР, vol.~153, issue~3, pp.~496-498, 1963.

%\bibitem{Ula2007}
%M.~Ulas,
%\emph{Rational Points on Certain Hyperelliptic Curves over Finite Fields},
%Bulletin of the Polish Academy of Sciences - Mathematics, vol.~55,
%issue~2, pp.~97-104, 2007.

%\bibitem{RistrettoDraft}
%H.~de Valence, J.~Grigg, M.~Hamburg, I.~Lovecruft, G.~Tankersly and
%F.~Valsorda,
%\emph{The ristretto255 and decaf448 Groups},\\
%\url{https://datatracker.ietf.org/doc/html/draft-irtf-cfrg-ristretto255-decaf448-03}

%\bibitem{Vel1971}
%J.~Vélu,
%\emph{Isogénies entre courbes elliptiques},
%C.R. Acad. Sc. Paris, Série A, vol.~273, pp.~238-241, 1971.

%\bibitem{Was2008}
%L.~Washington,
%\emph{Elliptic Curves: Number Theory and Cryptography} (2nd edition),
%Chapman and Hall/CRC, 2008.

%\bibitem{WhiWat1927}
%E.~Whittaker and G.~Watson,
%\emph{A Course of Modern Analysis},
%Cambridge University Press, 1927.

\bibitem{Wycheproof}
\emph{Project Wycheproof},\\
\url{https://github.com/C2SP/wycheproof}

\end{thebibliography}

% ----------------------------------------------------------------------

\end{document}
